\section{Root systems}\label{sec:3}
\begingroup
\emergencystretch=3em

We will construct the models before using the classification. All dimensions in this
section are dimensions of spans. We use both signs of every root, with squared norm 2,
as in Section 2. The symbol $\ell$ denotes rank, and $N$ denotes the number of lines.
Coordinates in the text start at 1; Lean's coordinates in \texttt{Fin} start at 0.
An isometric inclusion means a linear isometry of spans carrying the source roots
into the target roots. It is proper when that inclusion of root sets is proper.

\new\ We distinguish a proved blueprint result from an existing Lean declaration.
The tags \verb|\lean| give the Lean names; every one of them names a declaration in
Lean, and \path{dependency-diagram.md}, at the top of the repository, maps the numbered results
to their declarations. In
computational leaves we specify the integer lists, the Boolean test, and the soundness theorem.
The script \path{blueprint/checks/sec3/check_sec3.py} recomputes these tests exactly; it is a
sanity check, not a kernel proof of a new certificate.
No classification of Root Systems from Remark 3.11 is used here.

\begin{namedresult}{Lemma}{(transport on spans)}\label{lem:sec3-transport}
 \lean{LineSystem.Iso.spanIsometry}
 \uses{def:line-iso,def:line-system,def:indecomposable,def:2.1,def:2.2,hyp:2.4,facts:inner}
 \assume{$S,T$ are line systems in finite-dimensional real inner product spaces
 $V,V'$, and $f:\Phi(S)\to\Phi(T)$ is a bijection preserving inner products.}
 \prove{There is a unique surjective linear isometry
 $F:\operatorname{span}\Phi(S)\to\operatorname{span}\Phi(T)$ with $F(\Zr)=f(\Zr)$.
 It preserves line count and rank, star-closure, indecomposability, and maximality
 within the span. Its inverse has the same properties. For $\Za,\Zb\in\Phi(S)$ with
 $(\Za,\Zb)=-1$, the map $f$ carries the sets $\Ga$ and $\Delta$ of
 Hypothesis~\ref{hyp:2.4} for $\Za,\Zb$ bijectively onto those for $f(\Za),f(\Zb)$,
 and preserves orthogonality in $\Ga$.}
\end{namedresult}
\begin{proof}
 \pfsketch\ \new\ We state the transport fact separately so later results do not
 cite a construction hidden inside a proof. We will factor finite linear combinations through their kernels, and
 then transport the defining predicates.
 \step{s3trrelations}{For any real coefficients $a(\Zr)$, the two sums
 $\sum a(\Zr)\Zr$ and $\sum a(\Zr)f(\Zr)$ have equal squared norms.}
 \begin{proof}
  \pf\ Expand both squared norms by bilinearity from Section~\ref{facts:inner}; each term
  agrees by the inner-product hypothesis.~\qed
 \end{proof}
 \step{s3trmap}{The rule $F(\sum a(\Zr)\Zr)=\sum a(\Zr)f(\Zr)$ is a linear map of spans.}
 \begin{proof}
  \step{s3trkernel}{A zero source combination has zero target combination.}
  \begin{proof}\pf\ Apply zero squared norm from Section~\ref{facts:inner} to
   \stepref{s3trrelations}.~\qed\end{proof}
  \step{s3trquotient}{The target combination map factors through the quotient by
   the kernel of the source combination map.}
  \begin{proof}\pf\ Apply \path{Submodule.liftQ} to the kernel inclusion
   in \stepref{s3trkernel}.~\qed\end{proof}
  \qedstep
  \begin{proof}\pf\ Identify the quotient with the source range by
   \path{LinearMap.quotKerEquivRange}, and the range with the span by
   \path{Fintype.range_linearCombination}.~\qed\end{proof}
 \end{proof}
 \step{s3trisometry}{The map $F$ is an isometry and is onto the target span.}
 \begin{proof}
  \step{s3trinner}{It preserves the inner product of two combinations.}
  \begin{proof}\pf\ Bilinearity from Section~\ref{facts:inner} expands that product into
   the preserved products of pairs of roots.~\qed\end{proof}
  \step{s3tronto}{Each target combination has a source combination mapping to it.}
  \begin{proof}\pf\ Reindex its coefficients by the inverse of $f$.~\qed\end{proof}
  \qedstep
  \begin{proof}\pf\ Use \stepref{s3trinner} and \stepref{s3tronto}.~\qed\end{proof}
 \end{proof}
 \step{s3trpredicates}{The asserted invariants are preserved.}
 \begin{proof}
  \step{s3trlines}{Antipodal pairs correspond bijectively.}
  \begin{proof}\pf\ Linearity gives $F(-\Zr)=-F(\Zr)$; use the root bijection.~\qed\end{proof}
  \step{s3trrank}{The ranks agree.}
  \begin{proof}\pf\ A linear bijection carries a basis to a basis.~\qed\end{proof}
  \step{s3trclosed}{Star-closure is preserved in each direction.}
  \begin{proof}\pf\ Transport the sum and inner product in Definition~\ref{def:2.1}
   through $F$ or $F^{-1}$.~\qed\end{proof}
  \step{s3trindecomp}{Indecomposability is preserved in each direction.}
  \begin{proof}\pf\ Pull an orthogonal partition back through the corresponding
   root bijection in the indecomposability definition in Section~\ref{def:indecomposable}.~\qed\end{proof}
  \step{s3trmaximal}{Maximality within the span is preserved in each direction.}
  \begin{proof}\pf\ Transport a candidate vector of Definition~\ref{def:2.2}
   through the surjective span isometry.~\qed\end{proof}
  \step{s3trmarked}{\new\ The map $f$ carries $\Ga$ and $\Delta$ for $\Za,\Zb$ onto
   those for $f(\Za),f(\Zb)$, preserving orthogonality.}
  \begin{proof}\pf\ These sets are defined in Hypothesis~\ref{hyp:2.4} by inner products
   with $\Za$, $\Zb$ and $\Zc=-\Za-\Zb$; $f$ preserves inner products and
   $F(\Zc)=-f(\Za)-f(\Zb)$ by linearity.~\qed\end{proof}
  \qedstep
  \begin{proof}\pf\ The six preceding steps are the six assertions.~\qed\end{proof}
 \end{proof}
 \qedstep
 \begin{proof}\pf\ The finite-span lemma in Section~\ref{facts:inner} forces the formula
  for any linear extension; thus \stepref{s3trmap} is unique. Combine
  \stepref{s3trisometry} and \stepref{s3trpredicates}.~\qed\end{proof}
\end{proof}
\begingroup\raggedright
The quotient and combination declarations are in the Mathlib files
\path{LinearAlgebra/Quotient/Basic.lean}, \path{LinearAlgebra/Isomorphisms.lean},
and \path{LinearAlgebra/Finsupp/LinearCombination.lean}, respectively.\par
\endgroup

\begin{namedresult}{Lemma}{(restriction to the span)}\label{lem:sec3-restrict}
 \lean{LineSystem.restrictSpan}\lean{LineSystem.Iso.restrictSpan}
 \uses{def:line-system,lem:sec3-transport,facts:inner}
 \assume{$S$ is a line system in a finite-dimensional real inner product space $V$,
 and $W=\operatorname{span}\Phi(S)$, with the inner product of $V$.}
 \prove{The vectors of $\Phi(S)$, as vectors of $W$, form the root set of a line system
 $S|_W$ in $W$, and $\Phi(S|_W)$ spans $W$. The identity of $\Phi(S)$ is a bijection
 $\Phi(S|_W)\to\Phi(S)$ preserving inner products. Thus $S|_W$ and $S$ have the same
 number of lines, the same rank, and the same sets $\Ga$ and $\Delta$ with the same
 orthogonalities, and each is nonempty, star-closed or indecomposable exactly when the
 other is. A surjective linear isometry from $W$ that carries $\Phi(S|_W)$ onto a root
 set carries $\Phi(S)$ onto it.}
\end{namedresult}
\begin{proof}
 \pfsketch\ \new\ Theorem 2.13 assumes that the roots span the ambient space. We will
 state the passage to the span once, since Theorems 3.5 and 3.6 use it for $S$ and for
 the models.
 \step{s3reslines}{$S|_W$ is a line system in $W$.}
 \begin{proof}\pf\ The inner product of $W$ is that of $V$ on vectors of $W$, by
  \path{Submodule.coe_inner}; the conditions of the line-system definition in
  Section~\ref{def:line-system} are conditions on these inner products and on negation.~\qed\end{proof}
 \step{s3resspan}{$\Phi(S|_W)$ spans $W$.}
 \begin{proof}\pf\ This is \path{Submodule.span_span_coe_preimage} for the set
  $\Phi(S)$.~\qed\end{proof}
 \step{s3resinner}{The identity of $\Phi(S)$ preserves inner products.}
 \begin{proof}\pf\ Apply \path{Submodule.coe_inner} to two roots.~\qed\end{proof}
 \step{s3resinvariants}{The two systems have the asserted common invariants.}
 \begin{proof}\pf\ Apply Lemma~\ref{lem:sec3-transport} to the bijection of
  \stepref{s3resinner}; nonemptiness is preserved by any bijection.~\qed\end{proof}
 \step{s3resisometry}{The assertion on surjective linear isometries holds.}
 \begin{proof}\pf\ The two root sets are the same vectors of $W$, and
  $W=\operatorname{span}\Phi(S)$ by definition.~\qed\end{proof}
 \qedstep
 \begin{proof}\pf\ Steps \stepref{s3reslines}, \stepref{s3resspan},
  \stepref{s3resinner}, \stepref{s3resinvariants} and \stepref{s3resisometry} are the
  assertions.~\qed\end{proof}
\end{proof}
The two span declarations are in the Mathlib files
\path{Analysis/InnerProductSpace/Subspace.lean} and \path{LinearAlgebra/Span/Defs.lean}.

\begin{namedresult}{Lemma}{(an equal-rank proper inclusion)}\label{lem:sec3-proper-inclusion}
 \lean{LineSystem.not_isMaximal_of_proper_isometric_inclusion}
 \uses{def:line-system,def:2.2,facts:inner}
 \assume{\begin{pfenum}
  \item $S$ and $T$ are line systems in finite-dimensional real inner product spaces.
  \item The spans $W_S=\operatorname{span}\Phi(S)$ and $W_T=\operatorname{span}\Phi(T)$
   have the same dimension.
  \item $G:W_S\to W_T$ is a linear isometry with $G(\Phi(S))\subsetneq\Phi(T)$.
 \end{pfenum}}
 \prove{$S$ is not maximal within $W_S$.}
\end{namedresult}
\begin{proof}
 \pfsketch\ \new\ We will pull back a target root outside the image; it is a candidate
 of Definition 2.2 that is not a root of $S$.
 \step{s3propsurj}{$G$ is onto $W_T$.}
 \begin{proof}\pf\ An isometry is injective, and the dimensions agree by \asm[2]{0};
  apply \path{LinearMap.injective_iff_surjective_of_finrank_eq_finrank}.~\qed\end{proof}
 \step{s3proppick}{\pick\ $\Zt\in\Phi(T)\diff G(\Phi(S))$ and $\Zx\in W_S$ with
  $G(\Zx)=\Zt$.}
 \begin{proof}\pf\ The root $\Zt$ exists by \asm[3]{0}, and $\Zx$ by
  \stepref{s3propsurj}.~\qed\end{proof}
 \step{s3propnorm}{$(\Zx,\Zx)=2$.}
 \begin{proof}\pf\ $G$ preserves inner products and $(\Zt,\Zt)=2$.~\qed\end{proof}
 \step{s3propproducts}{For every $\Zr\in\Phi(S)$, $(\Zx,\Zr)\in\{-2,-1,0,1,2\}$.}
 \begin{proof}\pf\ $(\Zx,\Zr)=(\Zt,G(\Zr))$, and $\Zt,G(\Zr)\in\Phi(T)$ by
  \asm[3]{0}; use the line-system definition in Section~\ref{def:line-system}.~\qed\end{proof}
 \step{s3propnew}{$\Zx\notin\Phi(S)$.}
 \begin{proof}\pf\ Otherwise $\Zt=G(\Zx)$ would lie in $G(\Phi(S))$, contrary to the
  choice of $\Zt$.~\qed\end{proof}
 \qedstep
 \begin{proof}\pf\ Steps \stepref{s3propnorm}, \stepref{s3propproducts} and
  \stepref{s3propnew} give a vector of $W_S$ that Definition~\ref{def:2.2} requires in
  $\Phi(S)$ for a maximal $S$.~\qed\end{proof}
\end{proof}
The rank declaration is in the Mathlib file \path{LinearAlgebra/FiniteDimensional/Lemmas.lean}.

\begin{cgss}{Definition}{3.1}\label{def:3.1}
 \lean{LineSystem.An}\lean{LineSystem.Dn}\uses{def:line-system}
 \assume{$\ell$ is an integer, $\ell\geq1$ for $\typ A_\ell$ and $\ell\geq4$
 for $\typ D_\ell$; $(\Ze_i)$ is the standard orthonormal coordinate basis.}
 \pflet{The signed roots of the two systems are}
 \[
 \begin{split}
 \Phi(\typ A_\ell)&=\{\Ze_i-\Ze_j:1\leq i,j\leq\ell+1,\ i\ne j\},\\
 \Phi(\typ D_\ell)&=\{s\Ze_i+t\Ze_j:1\leq i<j\leq\ell,\ s,t\in\{1,-1\}\}.
 \end{split}
 \]
 These are line systems: each root has squared norm 2; taking negatives stays in
 the displayed set; pairwise products are integers of absolute value at most 2.
 We verify these assertions below with the other models.
 The inclusion $\Phi(\typ A_\ell)\subseteq\Phi(\typ D_{\ell+1})$ uses these coordinates.
 For $\ell<4$, the Lean definition \texttt{Dn} still uses this coordinate formula.
 The classification only uses $\typ D_\ell$ for $\ell\geq4$.
\end{cgss}

\begin{remark}
 \lean{LineSystem.cuboctahedronVertices}\lean{LineSystem.cuboctahedron}
 \lean{LineSystem.cuboctahedronDiagonals}
 \lean{LineSystem.cuboctahedronVertices_eq}\lean{LineSystem.ncard_cuboctahedronVertices}
 \lean{LineSystem.ncard_cubeEdges}\lean{LineSystem.IsCubeEdge.sym2_eq_of_midpoint_eq}
 \lean{LineSystem.cuboctahedron_Φ_inter_span_singleton}
 \lean{LineSystem.numLines_cuboctahedron}\lean{LineSystem.ncard_cuboctahedronDiagonals}
 \lean{LineSystem.cuboctahedron_eq_Dn}\lean{LineSystem.cuboctahedron.isoAn}
 \lean{LineSystem.nonempty_iso_cuboctahedron_An}
 The paper adds that $\typ A_3$ is represented by the 6 diagonals of the cuboctahedron. This is
 a side claim, and we prove it in \path{LineGraphs/SideClaims/Cuboctahedron.lean}. The vertices of
 the cuboctahedron are the midpoints of the twelve edges of the cube with vertices
 $(\pm1,\pm1,\pm1)$, and distinct edges have distinct midpoints. The midpoints are the twelve
 vectors with one coordinate 0 and two coordinates $\pm1$, which is $\Phi(\typ D_3)$ in the
 coordinates above. The vertex opposite to $\Zv$ is $-\Zv$, so each diagonal carries exactly two
 vertices, and the six diagonals are the lines of the line system $\typ D_3$
 (\path{LineSystem.cuboctahedron_eq_Dn}). This system is isomorphic to $\typ A_3$
 (\path{LineSystem.nonempty_iso_cuboctahedron_An}); in Lean the isomorphism is checked on the
 integer roots, by the kernel.
\end{remark}

\begin{cgss}{Definition}{3.2}\label{def:3.2}
 \lean{LineSystem.E₈viaA₈}\lean{LineSystem.E₈viaD₈}
 \lean{LineSystem.E₈viaSteiner}\lean{LineSystem.E₈}
 \lean{LineSystem.steiner348}\uses{def:3.1,def:line-system}
 \assume{Coordinate spaces have their usual real inner products.}
 \pflet{The following three signed sets, with the indicated normalization, define
 the three models of $\typ E_8$. The default is model (ii).}
 \begin{pfenum}
  \item In $\re^9$, take $\Phi(\typ A_8)$ together with both signs of
  $\frac13\sum_{i=1}^9\Ze_i-\sum_{i\in T}\Ze_i$, for $T\in\binom{[9]}3$.
  \item In $\re^8$, take $\Phi(\typ D_8)$ together with
  $\frac12\sum_{i=1}^8\epsilon(i)\Ze_i$, where $\epsilon(i)\in\{1,-1\}$ and
  $\prod_i\epsilon(i)=1$.
  \item In $\re^8$, take $\pm\sqrt2\Ze_i$ and
  $\frac1{\sqrt2}\sum_{i\in B}\epsilon(i)\Ze_i$ for all sign choices and all
  $B\in\mathcal B$, where the blocks of $\mathcal B=S(3,4,8)$ are
  \[
  \begin{array}{rrrrrrr}
  1234&1256&1278&1357&1368&1458&1467\\
  2358&2367&2457&2468&3456&3478&5678.
  \end{array}
  \]
 \end{pfenum}
 \repair\ This fixes the labelling required by Gap 5. We correct the second
 printed ``(i)'' to (ii) (Error 6), and remove the repeated sign condition
 (Error 13). \new\ The factor $1/\sqrt2$ in (iii) changes the printed vectors
 of squared norm 4 to our uniform squared norm 2, without changing their lines.
\end{cgss}

\begin{namedresult}{Lemma}{(integer realizations)}\label{lem:sec3-integer}
 \lean{Computations.C2.model_presentations}
 \uses{def:3.1,def:3.2,def:line-system,facts:inner}
 \assume{The coordinate sets are those of Definitions~\ref{def:3.1} and~\ref{def:3.2}.}
 \prove{They are line systems. Write $\Zr=c\Zv$ with integral $\Zv$ and
 $c^2q=1$. Their presentations are as follows; each list of raw vectors has no repetitions.
 \[
 \begin{array}{c|c|c|l}
 \text{model}&q&c&\text{raw vectors }\Zv\\\hline
 \typ A_\ell&1&1&\Ze_i-\Ze_j\ (i\ne j)\\
 \typ D_\ell&1&1&s\Ze_i+t\Ze_j\ (i<j)\\
 \typ E_8\text{ (i)}&9&1/3&3(\Ze_i-\Ze_j),\quad\pm(\mathbf1-3\mathbf1_T)\\
 \typ E_8\text{ (ii)}&4&1/2&2(s\Ze_i+t\Ze_j),\quad\epsilon\ (\prod\epsilon=1)\\
 \typ E_8\text{ (iii)}&2&1/\sqrt2&\pm2\Ze_i,\quad\epsilon\mathbf1_B
 \end{array}
 \]
 Every raw squared norm is $2q$, and every raw product lies in
 $q\{-2,-1,0,1,2\}$. Each triple of points belongs to exactly one block of $\mathcal B$;
 distinct blocks intersect in 0 or 2 points.}
\end{namedresult}
\begin{proof}
 \pfsketch\ \new\ We make the integer normalization and injective indices explicit
 to match the Lean models. We will check coordinates before using the common integer-model transfer.
 \step{s3intclassical}{The assertions hold for $\typ A_\ell,\typ D_\ell$.}
 \begin{proof}
  \step{s3intclassnorm}{Every such vector has exactly two nonzero coordinates, each a sign.}
  \begin{proof}\pf\ Read the two displays in Definition~\ref{def:3.1}.~\qed\end{proof}
  \step{s3intclassproducts}{Its squared norm is 2 and the products are integers between $-2$ and 2.}
  \begin{proof}\pf\ Sum the two coordinate products from \stepref{s3intclassnorm}.~\qed\end{proof}
  \step{s3intclassunique}{Supports and signs recover the displayed indices uniquely.}
  \begin{proof}\pf\ In $\typ A_\ell$ the positive and negative positions give the ordered
   pair; in $\typ D_\ell$ the sorted support gives the pair and its entries give the signs.~\qed\end{proof}
  \qedstep
  \begin{proof}\pf\ Negating the displayed signs stays in each list;
   use \stepref{s3intclassproducts} and \stepref{s3intclassunique}.~\qed\end{proof}
 \end{proof}
 \step{s3intone}{The assertions hold for model (i).}
 \begin{proof}
  \step{s3intonetriple}{For $t_T=\mathbf1-3\mathbf1_T$,
  $t_T\cdot t_U=9(|T\cap U|-1)$.}
  \begin{proof}\pf\ Expand the indicators and use $|T|=|U|=3$ in a 9-set.~\qed\end{proof}
  \step{s3intonemixed}{$3(\Ze_i-\Ze_j)\cdot t_T=9(\mathbf1_T(j)-\mathbf1_T(i))$.}
  \begin{proof}\pf\ Substitute the $i,j$ entries of $t_T$.~\qed\end{proof}
  \step{s3intoneunique}{The signed triple vectors are distinct and disjoint from the pair vectors.}
  \begin{proof}\pf\ A triple vector has three entries $-2$ and six entries 1;
   its negative has different entry values, while a pair vector has seven zero entries.~\qed\end{proof}
  \qedstep
  \begin{proof}\pf\ Use \stepref{s3intonetriple}, \stepref{s3intonemixed},
   \stepref{s3intoneunique}, and \stepref{s3intclassical} scaled by 3.~\qed\end{proof}
 \end{proof}
 \step{s3inttwo}{The assertions hold for model (ii).}
 \begin{proof}
  \step{s3inttwohalf}{For two parity-even sign vectors $u,v$, their Hamming distance
   $d$ is even, and $u\cdot v=8-2d\in\{-8,-4,0,4,8\}$.}
  \begin{proof}\pf\ The product of the eight coordinatewise products is 1,
   so the number of $-1$ products is even.~\qed\end{proof}
  \step{s3inttwomixed}{$2(s\Ze_i+t\Ze_j)\cdot\epsilon=2(s\epsilon(i)+t\epsilon(j))\in\{-4,0,4\}$.}
  \begin{proof}\pf\ Each summand in parentheses is a sign.~\qed\end{proof}
  \step{s3inttwounique}{The pair and half-root lists are disjoint and individually injective.}
  \begin{proof}\pf\ Pair roots have six zeros; half roots have none and their
   first seven signs determine the last by parity.~\qed\end{proof}
  \qedstep
  \begin{proof}\pf\ Use the last three steps and \stepref{s3intclassical};
   reversing all eight signs preserves parity.~\qed\end{proof}
 \end{proof}
 \step{s3intthree}{The assertions hold for model (iii).}
 \begin{proof}
  \step{s3intblocks}{The stated block sizes, intersections, and unique-triple property hold.}
  \begin{proof}\pf\ Enumerate the 14 displayed blocks and the 56 triples.
   These are the existing kernel results \path{LineSystem.steinerBlock_card},
   \path{LineSystem.steinerBlock_inter_even}, and
   \path{LineSystem.existsUnique_subset_steiner348} in \path{Models.lean}.~\qed\end{proof}
  \step{s3intthreeproducts}{Raw squared norms are 4, and products are even integers of absolute value at most 4.}
  \begin{proof}
   \step{s3intaxes}{Two axes have product 0 or $\pm4$; an axis and a block vector have product 0 or $\pm2$.}
   \begin{proof}\pf\ The axis has its only nonzero entry equal to $\pm2$.~\qed\end{proof}
   \step{s3intblockproducts}{Two block vectors have an even product between $-4$ and 4.}
   \begin{proof}\pf\ By \stepref{s3intblocks}, the product sums 0, 2, or 4 signs.~\qed\end{proof}
   \qedstep
   \begin{proof}\pf\ The two root types are exhausted by \stepref{s3intaxes}
    and \stepref{s3intblockproducts}; each self-product is 4.~\qed\end{proof}
  \end{proof}
  \step{s3intthreeunique}{The raw list is injective and closed under negation.}
  \begin{proof}\pf\ Support size distinguishes axes from blocks; support and entries
   recover every index and sign. All signs are allowed.~\qed\end{proof}
  \qedstep
  \begin{proof}\pf\ Use \stepref{s3intthreeproducts} and \stepref{s3intthreeunique}.~\qed\end{proof}
 \end{proof}
 \qedstep
 \begin{proof}\pf\ Apply \path{LineSystem.ofIntegerRoots} and
  \path{LineSystem.inner_integerVector} to the five checked presentations.
  These declarations are in \path{LineGraphs/LineSystem/IntegerModel.lean}.~\qed\end{proof}
\end{proof}

\begin{cgss}{Definition}{3.3}\label{def:3.3}
 \lean{LineSystem.E₇}\lean{LineSystem.E₆}
 \lean{LineSystem.T₇}\lean{LineSystem.T₆}
 \uses{def:3.2,lem:sec3-integer,def:line-system}
 \assume{$\typ E_8$ is model (ii) in $\re^8$.}
 \pflet{$\Zr_0=\Ze_7-\Ze_8$, $\Zs_0=\Ze_6-\Ze_7$, and}
 \[
 \Phi(\typ E_7)=\{\Zx\in\Phi(\typ E_8):x_7=x_8\},\qquad
 \Phi(\typ E_6)=\{\Zx\in\Phi(\typ E_8):x_6=x_7=x_8\}.
 \]
 Thus $\typ E_7=\typ E_8\cap\Zr_0^\perp$ and
 $\typ E_6=\typ E_8\cap\Zr_0^\perp\cap\Zs_0^\perp$.
 The roots $\Zr_0,\Zs_0,-\Zr_0-\Zs_0$ form a star.
 \repair\ We fix these choices first, and prove independence of both choices below
 (Gap 5); no transitivity hypothesis is part of the definition.
 For the printed counting argument we also name two cuts in model (iii):
 $T_7=\typ E_8^{\rm(iii)}\cap\Ze_8^\perp$ and
 $T_6=T_7\cap(\Ze_5+\Ze_6+\Ze_7+\Ze_8)^\perp$.
 Orthogonal cuts are line systems by restriction of the defining conditions.
\end{cgss}

\begin{namedresult}{Lemma}{(spans of the models)}\label{lem:sec3-spans}
 \lean{LineSystem.finrank_span_An}\lean{LineSystem.finrank_span_Dn}
 \lean{LineSystem.finrank_span_E₈}\lean{LineSystem.finrank_span_E₈viaD₈}
 \lean{LineSystem.finrank_span_E₈viaA₈}\lean{LineSystem.finrank_span_E₈viaSteiner}
 \lean{LineSystem.finrank_span_E₇}\lean{LineSystem.finrank_span_E₆}
 \uses{def:3.1,def:3.2,def:3.3,facts:inner}
 \assume{The models have the parameter ranges in Definitions~\ref{def:3.1}--\ref{def:3.3}.}
 \prove{Their spans and ranks are
 \[
 \begin{array}{c|c|c}
 \typ A_\ell&\{\Zx\in\re^{\ell+1}:\sum x_i=0\}&\ell\\
 \typ D_\ell&\re^\ell&\ell\\
 \typ E_8\text{ (i)}&\{\Zx\in\re^9:\sum x_i=0\}&8\\
 \typ E_8\text{ (ii), (iii)}&\re^8&8\\
 \typ E_7&\{\Zx\in\re^8:x_7=x_8\}&7\\
 \typ E_6&\{\Zx\in\re^8:x_6=x_7=x_8\}&6.
 \end{array}
 \]}
\end{namedresult}
\begin{proof}
 \pfsketch\ \repair\ We supply the dimension inputs deferred by Theorem 2.13
 (Gap 4). We will give spanning vectors with explicit independent coordinates.
 \step{s3spanA}{The $\typ A_\ell$ span is the stated hyperplane of dimension $\ell$.}
 \begin{proof}
  \step{s3spanAzero}{Every root has coordinate sum zero.}
  \begin{proof}\pf\ Its two nonzero coordinates are 1 and $-1$.~\qed\end{proof}
  \step{s3spanAbasis}{The vectors $\Ze_i-\Ze_{\ell+1}$, $1\leq i\leq\ell$, form a basis of this hyperplane.}
  \begin{proof}\pf\ A vector in the hyperplane is
   $\sum_{i=1}^\ell x_i(\Ze_i-\Ze_{\ell+1})$; its first $\ell$ coordinates make
   these coefficients unique.~\qed\end{proof}
  \qedstep
  \begin{proof}\pf\ The basis roots of \stepref{s3spanAbasis} and the containment
   in \stepref{s3spanAzero} give equality.~\qed\end{proof}
 \end{proof}
 \step{s3spanD}{The $\typ D_\ell$ roots span $\re^\ell$.}
 \begin{proof}\pf\ For $j\ne i$,
  $\Ze_i=((\Ze_i+\Ze_j)+(\Ze_i-\Ze_j))/2$ lies in the root span.~\qed\end{proof}
 \step{s3spanEi}{Model (i) spans the zero-sum hyperplane in $\re^9$.}
 \begin{proof}\pf\ Its additional roots have sum $9/3-3=0$, and it contains the
  $\typ A_8$ hyperplane of \stepref{s3spanA}.~\qed\end{proof}
 \step{s3spanEii}{Model (ii) spans $\re^8$.}
 \begin{proof}\pf\ It contains $\typ D_8$, whose span is given by \stepref{s3spanD}.~\qed\end{proof}
 \step{s3spanEiii}{Model (iii) spans $\re^8$.}
 \begin{proof}\pf\ Its eight axis roots are nonzero scalar multiples of a coordinate basis.~\qed\end{proof}
 \step{s3spanSeven}{The $\typ E_7$ span has basis $\Ze_1,\ldots,\Ze_6,\Ze_7+\Ze_8$.}
 \begin{proof}
  \step{s3spanSevenin}{These seven vectors lie in the span.}
  \begin{proof}\pf\ The $\typ D_6$ roots on the first six coordinates give the six axes
   by \stepref{s3spanD}; $\Ze_7+\Ze_8$ is itself an $\typ E_7$ root.~\qed\end{proof}
  \qedstep
  \begin{proof}\pf\ The displayed vectors uniquely expand every vector with
   $x_7=x_8$, which contains all the roots by Definition~\ref{def:3.3}.~\qed\end{proof}
 \end{proof}
 \step{s3spanSix}{The $\typ E_6$ span has basis $\Ze_1,\ldots,\Ze_5,\Ze_6+\Ze_7+\Ze_8$.}
 \begin{proof}
  \step{s3spanSixaxes}{The first five axes belong to the span.}
  \begin{proof}\pf\ Apply \stepref{s3spanD} to the $\typ D_5$ roots on the first five coordinates.~\qed\end{proof}
  \step{s3spanSixtail}{The last basis vector belongs to the span.}
  \begin{proof}\pf\ Subtract $\sum_{i=1}^5\Ze_i$ from twice the root
   $\frac12\mathbf1$, using \stepref{s3spanSixaxes}.~\qed\end{proof}
  \qedstep
  \begin{proof}\pf\ These six vectors uniquely expand every vector with
   $x_6=x_7=x_8$, which contains the roots by Definition~\ref{def:3.3}.~\qed\end{proof}
 \end{proof}
 \qedstep
 \begin{proof}\pf\ The six cases exhaust the models. The cited dimension declarations
  already implement these arguments in \path{LineGraphs/LineSystem/ModelSpans.lean}.~\qed\end{proof}
\end{proof}

\begin{namedresult}{Computation}{C2}\label{comp:C2}
 \lean{Computations.C2.line_counts}
 \lean{LineSystem.numLines_T₇}\lean{LineSystem.numLines_T₆}
 \uses{def:3.1,def:3.2,def:3.3,lem:sec3-integer}
 \assume{The models and the fixed cuts $T_7,T_6$ are as defined above.}
 \prove{Their line counts are $\ell(\ell+1)/2$, $\ell(\ell-1)$,
 120 in each $\typ E_8$ model, 63 for $\typ E_7,T_7$, and 36 for $\typ E_6,T_6$.}
\end{namedresult}
\begin{proof}
 \pfsketch\ We will count injective signed indices and divide by 2.
 \step{s3countA}{$\typ A_\ell$ has $\ell(\ell+1)/2$ lines.}
 \begin{proof}\pf\ There are $(\ell+1)\ell$ ordered distinct pairs in
  Definition~\ref{def:3.1}; injectivity is Lemma~\ref{lem:sec3-integer}.~\qed\end{proof}
 \step{s3countD}{$\typ D_\ell$ has $\ell(\ell-1)$ lines.}
 \begin{proof}\pf\ There are $4\binom\ell2$ signed indices; divide by 2.~\qed\end{proof}
 \step{s3countEone}{Model (i) has $36+84=120$ lines.}
 \begin{proof}\pf\ Its disjoint types contribute $\binom92$ and $\binom93$ antipodal pairs.~\qed\end{proof}
 \step{s3countEtwo}{Model (ii) has $56+64=120$ lines.}
 \begin{proof}\pf\ Its disjoint types contribute $2\binom82$ and $2^7/2$ antipodal pairs.~\qed\end{proof}
 \step{s3countEthree}{Model (iii) has $8+8\cdot14=120$ lines.}
 \begin{proof}\pf\ Each of the 14 supports contributes $2^4/2$ pairs, in addition to the eight axes.~\qed\end{proof}
 \step{s3countSeven}{$\typ E_7$ has $(60+2+64)/2=63$ lines.}
 \begin{proof}
  \step{s3countSevenpairs}{Its pair roots are $\typ D_6$ and $\pm(\Ze_7+\Ze_8)$.}
  \begin{proof}\pf\ The condition $x_7=x_8$ permits a support inside $[6]$, or the
   support $\{7,8\}$ with equal signs.~\qed\end{proof}
  \step{s3countSevenhalf}{It has $2^5\cdot2=64$ half roots.}
  \begin{proof}\pf\ The first six signs have even parity; the last two are equal with two choices.~\qed\end{proof}
  \qedstep
  \begin{proof}\pf\ Add the two disjoint counts and divide by 2.~\qed\end{proof}
 \end{proof}
 \step{s3countSix}{$\typ E_6$ has $(40+32)/2=36$ lines.}
 \begin{proof}
  \step{s3countSixpairs}{Its pair roots are exactly $\typ D_5$ on $[5]$.}
  \begin{proof}\pf\ A two-element support cannot meet the three required equal
   nonzero coordinates $6,7,8$.~\qed\end{proof}
  \step{s3countSixhalf}{There are $2^5=32$ half roots.}
  \begin{proof}\pf\ The first five signs determine the common last sign by the even-parity condition.~\qed\end{proof}
  \qedstep
  \begin{proof}\pf\ Add the disjoint counts and divide by 2.~\qed\end{proof}
 \end{proof}
 \step{s3countTseven}{$T_7$ has $7+8\cdot7=63$ lines.}
 \begin{proof}\pf\ Exactly seven blocks in Definition~\ref{def:3.2} avoid 8;
  each contributes eight lines, in addition to seven axes.~\qed\end{proof}
 \step{s3countTsix}{$T_6$ has $4+8+4\cdot6=36$ lines.}
 \begin{proof}
  \step{s3countTsixaxes}{The four surviving axes have indices $1,2,3,4$.}
  \begin{proof}\pf\ Impose $x_8=0$ and $x_5+x_6+x_7=0$ on an axis.~\qed\end{proof}
  \step{s3countTsixblock}{The block 1234 contributes eight lines.}
  \begin{proof}\pf\ Both linear conditions vanish on every sign choice on this block.~\qed\end{proof}
  \step{s3countTsixother}{Each of the other six blocks avoiding 8 contributes four lines.}
  \begin{proof}\pf\ Each meets $\{5,6,7\}$ twice; those signs must be opposite.
   There are $2\cdot2^2=8$ signed roots per block.~\qed\end{proof}
  \qedstep
  \begin{proof}\pf\ The three disjoint root types give the stated sum.~\qed\end{proof}
 \end{proof}
 \qedstep
 \begin{proof}\pf\ All counts use the injectivity in Lemma~\ref{lem:sec3-integer}.
  The common Lean transfer is \path{LineSystem.IntegerModel.numLines_eq}.
  The finite cardinal certificates \leancite{LineSystem.card_E₇Index}{\texttt{LineSystem.\allowbreak card\_\allowbreak E\ensuremath{{}_7}Index}} and
  \leancite{LineSystem.card_E₆Index}{\texttt{LineSystem.\allowbreak card\_\allowbreak E\ensuremath{{}_6}Index}} already check 126 and 72 in the kernel.~\qed\end{proof}
\end{proof}

\begin{namedresult}{Computation}{C3}\label{comp:C3}
 \lean{LineSystem.E₈viaA₈.isoE₈}\lean{LineSystem.E₈viaSteiner.isoE₈}
 \uses{def:3.2,lem:sec3-integer,lem:sec3-spans,comp:C2,lem:sec3-transport}
 \assume{The three models have the normalization of Definition~\ref{def:3.2}.}
 \prove{Model (i) is isometric to model (ii) by
 \[
 U(\Zx)_i=d_i(x_i+x_9/2)\quad(1\leq i\leq8),\qquad
 d_1=\cdots=d_7=1,\quad d_8=-1,
 \]
 on the zero-sum hyperplane. Model (iii) is isometric to model (ii) by
 $Q/\sqrt2$, where $Q$ is block diagonal with four blocks
 $\left(\begin{smallmatrix}1&1\\1&-1\end{smallmatrix}\right)$
 on the coordinate pairs $12,34,56,78$. Both maps carry the full signed root sets onto one another.}
\end{namedresult}
\begin{proof}
 \pfsketch\ \new\ We give explicit maps, so equivalence of the definitions does not
 use Theorem 3.5. In raw coordinates their formulas are rational, even though the
 normalization of the Steiner map introduces $\sqrt2$.
 \step{s3c3Umetric}{$U$ preserves inner products on the zero-sum hyperplane.}
 \begin{proof}\pf\ For $\sum x_i=\sum y_i=0$, expansion gives
  $\sum_{i=1}^8(x_i+x_9/2)(y_i+y_9/2)=\sum_{i=1}^9x_iy_i$.
  The factors $d_i^2=1$ do not change this sum.~\qed\end{proof}
 \step{s3c3Uroots}{$U$ sends every root of model (i) to a root of model (ii).}
 \begin{proof}
  \step{s3c3Upair}{\case{$\Zx=\Ze_i-\Ze_j$, with $i,j\leq8$.}}
  \begin{proof}\pf\ Its image is $d_i\Ze_i-d_j\Ze_j$, a $\typ D_8$ root.~\qed\end{proof}
  \step{s3c3Unine}{\case{$\Zx=\Ze_i-\Ze_9$, with $i\leq8$.}}
  \begin{proof}\pf\ Before the factor $d$, the image has one entry $1/2$ and
   seven entries $-1/2$. Negating coordinate 8 makes its sign parity even.~\qed\end{proof}
  \step{s3c3Utripleout}{\case{$\Zx=\mathbf1/3-\mathbf1_T$ and $9\notin T$.}}
  \begin{proof}\pf\ Before $d$, the image is $\mathbf1/2-\mathbf1_T$, with
   three negative signs; the coordinate-8 sign change makes their number even.~\qed\end{proof}
  \step{s3c3Utriplein}{\case{$\Zx=\mathbf1/3-\mathbf1_T$ and $T=\{i,j,9\}$.}}
  \begin{proof}\pf\ Its image is $-d_i\Ze_i-d_j\Ze_j$, a $\typ D_8$ root.~\qed\end{proof}
  \step{s3c3Usigns}{Negatives of the listed roots also map to roots.}
  \begin{proof}\pf\ $U$ is linear and model (ii) is closed under negation by
   Lemma~\ref{lem:sec3-integer}.~\qed\end{proof}
  \qedstep
  \begin{proof}\pf\ The pair and triple alternatives in Definition~\ref{def:3.2}
   are exhausted by the five preceding steps.~\qed\end{proof}
 \end{proof}
 \step{s3c3Qmetric}{$Q/\sqrt2$ preserves inner products.}
 \begin{proof}\pf\ Each displayed block has transpose times itself equal to $2I_2$,
  so $Q^TQ=2I_8$.~\qed\end{proof}
 \step{s3c3Qroots}{$Q/\sqrt2$ sends the roots of model (iii) into model (ii).}
 \begin{proof}
  \step{s3c3Qaxis}{An axis root $\pm\sqrt2\Ze_i$ maps to a $\typ D_8$ root.}
  \begin{proof}\pf\ The $i$th column of $Q$ has two entries in $\{1,-1\}$.~\qed\end{proof}
  \step{s3c3Qunion}{\case{$B$ is the union of two pairs among $12,34,56,78$.}}
  \begin{proof}\pf\ On each occupied pair $Q$ sends two signs to one entry $\pm2$
   and one zero; division by 2 gives a $\typ D_8$ root.~\qed\end{proof}
  \step{s3c3Qtransversal}{\case{$B$ meets each of the four pairs once.}}
  \begin{proof}
   \step{s3c3Qeven}{The number of second positions selected by $B$ is even.}
   \begin{proof}\pf\ Read the eight transversal blocks in Definition~\ref{def:3.2}:
    \[\begin{array}{rrrr}1357&1368&1458&1467\\2358&2367&2457&2468.\end{array}\]
    \qed\end{proof}
   \qedstep
   \begin{proof}\pf\ Selecting a first position gives a pair of signs with product 1;
    selecting a second gives product $-1$. Step \stepref{s3c3Qeven} makes the total
    product 1, so $Q\Zv/2$ is a half root of model (ii).~\qed\end{proof}
  \end{proof}
  \qedstep
  \begin{proof}\pf\ The other six blocks are $1234,1256,1278,3456,3478,5678$,
   the unions in \stepref{s3c3Qunion}; together with axes these exhaust the raw types.~\qed\end{proof}
 \end{proof}
 \step{s3c3onto}{Both root maps are onto.}
 \begin{proof}\pf\ Each is injective by its isometry property and maps between two
  finite sets of 240 roots by C2, so each is surjective.~\qed\end{proof}
 \qedstep
 \begin{proof}\pf\ Use \stepref{s3c3onto} and the span descriptions in
  Lemma~\ref{lem:sec3-spans}. Equivalently, the induced index bijections satisfy
  the hypothesis of the existing constructor \path{LineSystem.Iso.ofIntegerModels}.~\qed\end{proof}
\end{proof}
For a kernel certificate the raw maps are
$\Zv\mapsto(d_i(2v_i+v_9)/3)_{i=1}^8$ and $\Zv\mapsto Q\Zv$.
The 240 images in each case are integral. In Lean (\path{Computations/C3.lean}) the kernel
checks that the lists of images are permutations of the roots of model (ii), and checks the
scaled Gram equalities; the script \path{blueprint/checks/sec3/check_sec3.py} checks the images,
their bijections, and the scaled Gram equalities in exact arithmetic.

\begin{namedresult}{Lemma}{(a connected spanning set)}\label{lem:sec3-connected-span}
 \lean{LineSystem.isIndecomposable_of_connected_spanning}
 \uses{def:line-system,def:indecomposable,facts:inner}
 \assume{$S$ is a line system in a finite-dimensional real inner product space;
 $B\subseteq\Phi(S)$ is nonempty, spans $\operatorname{span}\Phi(S)$, and its
 graph joining distinct roots of nonzero inner product is connected.}
 \prove{$S$ is indecomposable.}
\end{namedresult}
\begin{proof}
 \pfsketch\ \new\ We isolate the spanning criterion for the model verifications.
 An orthogonal partition cannot separate an edge of $B$, or a root
 from every vector of a spanning set.
 \step{s3connpartition}{\assume{$P\cup Q=\Phi(S)$ and $P\perp Q$.}
  \prove{$P=\varnothing$ or $Q=\varnothing$.}}
 \begin{proof}
  \step{s3connunique}{No root belongs to both parts.}
  \begin{proof}\pf\ Such a root would have self-product both 0 and 2.~\qed\end{proof}
  \step{s3connedge}{The endpoints of every edge of $B$ lie in the same part.}
  \begin{proof}\pf\ Different parts have inner product zero by \asm{1}.~\qed\end{proof}
  \step{s3connpath}{All of $B$ lies in one part.}
  \begin{proof}\pf\ Choose a vertex of the nonempty $B$ and propagate its part
   along a path to each vertex, using \stepref{s3connedge}.~\qed\end{proof}
  \step{s3connother}{The other part is empty.}
  \begin{proof}
   \step{s3connzero}{Every vector of that part is orthogonal to the whole span.}
   \begin{proof}\pf\ It is orthogonal to $B$ by \asm{1} and
    \stepref{s3connpath}; extend by the finite-span lemma in Section~\ref{facts:inner}.~\qed\end{proof}
   \qedstep
   \begin{proof}\pf\ A root in that part is in the span, so
    \stepref{s3connzero} would give its self-product as 0 instead of 2.~\qed\end{proof}
  \end{proof}
  \qedstep
  \begin{proof}\pf\ Use \stepref{s3connother}, whichever part contains $B$.~\qed\end{proof}
 \end{proof}
 \qedstep
 \begin{proof}\pf\ Step \stepref{s3connpartition} is the indecomposability definition in Section~\ref{def:indecomposable}.~\qed\end{proof}
\end{proof}

\begin{namedresult}{Computation}{C4}\label{comp:C4}
 \lean{Computations.C4.exceptional_star_data}
 \lean{Computations.C4.marked_star_data}
 \lean{LineSystem.E₆.latticeGraphIso}\lean{LineSystem.E₇.triangularGraphComplIso}
 \lean{LineSystem.Hyp24.nonempty_latticeGraph_iso_of_S_eq_E₆}
 \lean{LineSystem.Hyp24.nonempty_triangularGraph_compl_iso_of_S_eq_E₇}
 \lean{LineSystem.E₆.exists_isometry_star}\lean{LineSystem.E₇.exists_isometry_star}
 \lean{LineSystem.E₈.exists_isometry_star}
 \lean{LineSystem.Hyp24.orthGraph_isSRGWith_of_S_eq_E₆}\lean{LineSystem.Hyp24.orthGraph_isSRGWith_of_S_eq_E₇}
 \lean{LineSystem.Hyp24.orthGraph_isSRGWith_of_S_eq_E₈}
 \lean{LineSystem.Hyp24.card_Γa_of_S_eq_E₆}\lean{LineSystem.Hyp24.card_Γa_of_S_eq_E₇}
 \lean{LineSystem.Hyp24.card_Γa_of_S_eq_E₈}
 \lean{LineSystem.Hyp24.card_Δ_of_S_eq_E₆}\lean{LineSystem.Hyp24.card_Δ_of_S_eq_E₇}
 \lean{LineSystem.Hyp24.card_Δ_of_S_eq_E₈}
 \lean{LineSystem.isStarClosed_An}\lean{LineSystem.isStarClosed_Dn}
 \lean{LineSystem.isStarClosed_E₈}\lean{LineSystem.isStarClosed_E₇}\lean{LineSystem.isStarClosed_E₆}
 \lean{LineSystem.isStarClosed_E₈viaA₈}\lean{LineSystem.isStarClosed_E₈viaSteiner}
 \lean{LineSystem.isIndecomposable_An}\lean{LineSystem.isIndecomposable_Dn}
 \lean{LineSystem.isIndecomposable_E₈}\lean{LineSystem.isIndecomposable_E₇}
 \lean{LineSystem.isIndecomposable_E₆}\lean{LineSystem.isIndecomposable_E₈viaA₈}
 \lean{LineSystem.isIndecomposable_E₈viaSteiner}
 \lean{LineSystem.An.hyp24}\lean{LineSystem.Dn.hyp24}\lean{LineSystem.E₆.hyp24}
 \lean{LineSystem.E₇.hyp24}\lean{LineSystem.E₈.hyp24}
 \lean{LineSystem.An.orthGraphIso}\lean{LineSystem.Dn.orthGraphIso}
 \lean{LineSystem.E₆.orthGraphIso}\lean{LineSystem.E₇.orthGraphIso}\lean{LineSystem.E₈.orthGraphIso}
 \lean{LineSystem.An.card_Γa}\lean{LineSystem.Dn.card_Γa}\lean{LineSystem.E₆.card_Γa}
 \lean{LineSystem.E₇.card_Γa}\lean{LineSystem.E₈.card_Γa}
 \lean{LineSystem.E₆.card_Δ_div_two}\lean{LineSystem.E₇.card_Δ_div_two}
 \lean{LineSystem.E₈.card_Δ_div_two}\lean{LineSystem.E₆.orthGraph_isSRGWith}
 \lean{LineSystem.E₇.orthGraph_isSRGWith}\lean{LineSystem.E₈.orthGraph_isSRGWith}
 \uses{def:3.1,def:3.2,def:3.3,lem:sec3-integer,lem:sec3-spans,comp:C2,comp:C3,lem:sec3-transport,lem:sec3-connected-span,hyp:2.4,def:gamma-graph,def:2.1}
 \assume{We use the models in their specified ranges, and $X=\opk{Orth}(\Ga)$
 has the convention of Hypothesis~\ref{hyp:2.4}: distinct vertices are adjacent when
 their inner product is 0.}
 \prove{Every model is nonempty, star-closed, and indecomposable. The model
 $\typ A_1$ has one line and rank 1. For the marked star
 $\Za=\Ze_1-\Ze_2$, $\Zb=\Ze_2-\Ze_3$ the following data hold:
 \[
 \begin{array}{c|c|c|c|c}
 \text{model}&\text{rank}&\Ga\text{ or its graph}&|\Ga|&|\Delta|/2\\\hline
 \typ A_\ell\ (\ell\geq2)&\ell&\{\Ze_j-\Ze_3:4\leq j\leq\ell+1\}&\ell-2&\text{not needed}\\
 \typ D_\ell&\ell&\{\Ze_1+\Ze_2\}\cup\{-\Ze_3\pm\Ze_j:4\leq j\leq\ell\}&2\ell-5&\text{not needed}\\
 \typ E_6&6&\text{strongly regular }(9,4,1,2)&9&6\\
 \typ E_7&7&\text{strongly regular }(15,6,1,3)&15&15\\
 \typ E_8\text{ (ii)}&8&\text{strongly regular }(27,10,1,5)&27&36
 \end{array}
 \]
 The $\typ A$ graph is empty and the $\typ D$ graph is the friendship graph of $\ell-3$ triangles.
 The three exceptional rows hold for \emph{every} ordered pair $\Za,\Zb$ of roots
 with inner product $-1$ in the corresponding fixed model. The signed permutation
 $\sigma$ with $\sigma(\Ze_1)=\Ze_2$, $\sigma(\Ze_2)=-\Ze_1$, $\sigma(\Ze_3)=-\Ze_3$,
 fixing the other coordinate vectors, is an automorphism of $\typ D_\ell$, $\typ E_8$,
 $\typ E_7$ and $\typ E_6$ that carries $\Za,\Zb$ to the ordered star
 $\Za'=\Ze_1+\Ze_2$, $\Zb'=-\Ze_1+\Ze_3$. For this star, the graph $\opk{Orth}(\Ga)$ of
 $\typ E_6$ is the $3\times3$ lattice graph (adjacency in one common coordinate), and that
 of $\typ E_7$ is the disjointness graph on $\binom{[6]}2$, the complement of $T(6)$. Finally,
 the isometries of each exceptional model are transitive on its ordered stars: for
 $t\in\{6,7,8\}$ and $\Zv,\Zw\in\Phi(\typ E_t)$ with $(\Zv,\Zw)=-1$ there is a linear isometry
 $g$ of $\re^8$ with $g(\Phi(\typ E_t))=\Phi(\typ E_t)$, $g(\Za')=\Zv$ and $g(\Zb')=\Zw$.}
\end{namedresult}
\begin{proof}
 \pfsketch\ We will prove the classical assertions parametrically. Exceptional
 star data are finite tests on the complete integer root lists. Indecomposability
 has a short spanning-set proof, independent of any classification.
 \step{s3c4closed}{The five default models are star-closed.}
 \begin{proof}
  \step{s3c4Dclosed}{$\typ D_\ell$ is star-closed.}
  \begin{proof}
   \step{s3c4Dnorm}{\assume{$u,v$ are integral norm-2 roots with $u\cdot v=-1$.}
    \prove{$(u+v)\cdot(u+v)=2$.}}
   \begin{proof}\pf\ Bilinearity gives $2+2(-1)+2=2$.~\qed\end{proof}
   \qedstep
   \begin{proof}\pf\ An integral norm-2 vector has exactly two nonzero entries,
    each $\pm1$: this is \path{LineSystem.exists_dnRoot_of_inner_self}.
    Apply it to \stepref{s3c4Dnorm}.~\qed\end{proof}
  \end{proof}
  \step{s3c4Aclosed}{$\typ A_\ell$ is star-closed.}
  \begin{proof}\pf\ In \stepref{s3c4Dclosed} the sum also has coordinate sum zero,
   so its two signs are opposite. The Lean criterion is
   \path{LineSystem.exists_anRoot_of_inner_self_of_sum}.~\qed\end{proof}
  \step{s3c4Eclosed}{$\typ E_8$ is star-closed.}
  \begin{proof}\pf\ The existing certificate
   \leancite{Computations.C4.e₈_star_checked}{\texttt{Computations.\allowbreak C4.\allowbreak e\ensuremath{{}_8}\_\allowbreak star\_\allowbreak checked}} checks
   \path{Computations.checkStar} with dimension 8, target $-4$, the complete
   doubled root list, and its sound root decoder.
   Apply \path{LineSystem.IntegerModel.isStarClosed_of_checkStar}.~\qed\end{proof}
  \step{s3c4Sevenclosed}{$\typ E_7$ is star-closed.}
  \begin{proof}\pf\ Apply \path{LineSystem.IsStarClosed.orthogonalTo} to
   \stepref{s3c4Eclosed} and $\Zr_0$.~\qed\end{proof}
  \step{s3c4Sixclosed}{$\typ E_6$ is star-closed.}
  \begin{proof}\pf\ Apply the same named orthogonal-cut theorem to
   \stepref{s3c4Sevenclosed} and $\Zs_0$.~\qed\end{proof}
  \qedstep
  \begin{proof}\pf\ These five cases exhaust the default models.~\qed\end{proof}
 \end{proof}
 \step{s3c4indec}{The five default models are nonempty and indecomposable.}
 \begin{proof}
  \pflet{$B_m=\{\Ze_i-\Ze_m:1\leq i<m\}\cup\{\Ze_1+\Ze_m\}$ for $m\geq3$,
   and $\Zh=\frac12\mathbf1\in\re^8$.}
  \step{s3c4Aindec}{$\typ A_\ell$ has a nonempty connected spanning set.}
  \begin{proof}\pf\ The basis $\{\Ze_i-\Ze_{\ell+1}:1\leq i\leq\ell\}$ in
   Lemma~\ref{lem:sec3-spans} has product 1 between distinct members;
   for $\ell=1$ its graph has one vertex.~\qed\end{proof}
  \step{s3c4Bindec}{$B_m$ is a connected spanning set in $\typ D_m$.}
  \begin{proof}
   \step{s3c4Bconnect}{Its difference roots form a clique, and $\Ze_1+\Ze_m$
    has product $-1$ with $\Ze_2-\Ze_m$.}
   \begin{proof}\pf\ Compute the products from the coordinate basis, using $m\geq3$.~\qed\end{proof}
   \step{s3c4Bspan}{$B_m$ spans $\re^m$.}
   \begin{proof}\pf\ The difference of $\Ze_1+\Ze_m$ and $\Ze_1-\Ze_m$
    gives $2\Ze_m$; adding this recovered axis to each difference gives every other axis.~\qed\end{proof}
   \qedstep
   \begin{proof}\pf\ Use \stepref{s3c4Bconnect} and \stepref{s3c4Bspan}.~\qed\end{proof}
  \end{proof}
  \step{s3c4Eindec}{$B_8$ is a connected spanning subset of $\typ E_8$.}
  \begin{proof}\pf\ It lies in $\typ D_8\subseteq \typ E_8$ and has the properties in
   \stepref{s3c4Bindec}.~\qed\end{proof}
  \step{s3c4Sevenindec}{$B_6\cup\{\Ze_7+\Ze_8,\Zh\}$ is a connected spanning subset of $\typ E_7$.}
  \begin{proof}
   \step{s3c4Sevenjoin}{$\Zh$ has product 1 with both $\Ze_1+\Ze_6\in B_6$
    and $\Ze_7+\Ze_8$.}
   \begin{proof}\pf\ Each product is $1/2+1/2$.~\qed\end{proof}
   \qedstep
   \begin{proof}\pf\ Step \stepref{s3c4Bindec} spans the first six axes;
    adding $\Ze_7+\Ze_8$ spans the space in Lemma~\ref{lem:sec3-spans}.
    Step \stepref{s3c4Sevenjoin} connects the parts.~\qed\end{proof}
  \end{proof}
  \step{s3c4Sixindec}{$B_5\cup\{\Zh\}$ is a connected spanning subset of $\typ E_6$.}
  \begin{proof}\pf\ $B_5$ spans the first five axes by \stepref{s3c4Bindec};
   $\Zh$ has nonzero last three equal coordinates and has product 1 with
   $\Ze_1+\Ze_5$. Use the $\typ E_6$ span in Lemma~\ref{lem:sec3-spans}.~\qed\end{proof}
  \qedstep
  \begin{proof}\pf\ Apply Lemma~\ref{lem:sec3-connected-span} to each displayed nonempty set.~\qed\end{proof}
 \end{proof}
 \step{s3c4othermodels}{The other two $\typ E_8$ models are nonempty, star-closed, and indecomposable.}
 \begin{proof}\pf\ Transport \stepref{s3c4closed} and \stepref{s3c4indec}
  through C3, using Lemma~\ref{lem:sec3-transport}.~\qed\end{proof}
 \step{s3c4Agraph}{The marked $\typ A_\ell$ graph has the asserted vertices and no edges.}
 \begin{proof}
  \step{s3c4Avertices}{The equations $x_1=x_2$ and $x_2-x_3=1$ select precisely
   $\Ze_j-\Ze_3$ for $4\leq j\leq\ell+1$.}
  \begin{proof}\pf\ A difference root cannot have two equal nonzero entries at 1 and 2,
   so $x_1=x_2=0$ and $x_3=-1$.~\qed\end{proof}
  \qedstep
  \begin{proof}\pf\ Two distinct selected roots have product 1; their number is $\ell-2$.~\qed\end{proof}
 \end{proof}
 \step{s3c4Dgraph}{The marked $\typ D_\ell$ graph is the asserted friendship graph.}
 \begin{proof}
  \step{s3c4Dvertices}{The equations select $\Ze_1+\Ze_2$ and $-\Ze_3\pm\Ze_j$ for $j\geq4$.}
  \begin{proof}\pf\ Either $x_1=x_2\ne0$, forcing $x_1=x_2=1,x_3=0$;
   or $x_1=x_2=0$, forcing $x_3=-1$ and one further sign outside $[3]$.~\qed\end{proof}
  \step{s3c4Dcentre}{The vertex $\Ze_1+\Ze_2$ is adjacent to all the others.}
  \begin{proof}\pf\ Their supports are disjoint.~\qed\end{proof}
  \step{s3c4Dpairs}{Among the other vertices, adjacency pairs the two signs of the same $j$.}
  \begin{proof}\pf\ The product is 0 for those two signs, and 1 for different indices.~\qed\end{proof}
  \qedstep
  \begin{proof}\pf\ Steps \stepref{s3c4Dcentre} and \stepref{s3c4Dpairs} give
   $\ell-3$ triangles and $1+2(\ell-3)$ vertices.~\qed\end{proof}
 \end{proof}
 \step{s3c4exceptional}{\new\ Every ordered star in each default exceptional model
  has its stated row. We check all choices to avoid assuming transitivity.}
 \begin{proof}
  \pflet{$R$ be the sorted complete doubled integer root list of the model;
   for $a,b\in R$ with $a\cdot b=-4$, set
   $L(a,b)=\{v\in R:v\cdot a=0,\ v\cdot b=4\}$ and
   $D(a,b)=\{v\in R:v\cdot a=v\cdot b=0\}$.}
  A row test compares the two list lengths with $m,2\delta$; for each $u\in L$
  it counts $v\in L$ with $u\cdot v=0$ and compares with $d$; for each distinct
  $u,v\in L$ it counts $w\in L$ with $u\cdot w=v\cdot w=0$ and compares with
  $\lambda$ when $u\cdot v=0$, and with $\mu$ otherwise.
  \step{s3c4lists}{These lists represent exactly $\Ga$ and $\Delta$, respectively.}
  \begin{proof}\pf\ Divide raw dot products by $q=4$ using
   \path{LineSystem.inner_integerVector}; the third equation of $\Ga$ follows
   from $\Zc=-\Za-\Zb$. Completeness is the model's
   \texttt{mem\_rootList} theorem in \path{Computations/ModelTables.lean}.~\qed\end{proof}
  \step{s3c4census}{The complete finite tests give
  \[
  \begin{array}{c|r|r|r|r|r|r}
  R&\#(a,b)&|L|&|D|&d&\lambda&\mu\\\hline
  \typ E_6&1440&9&12&4&1&2\\
  \typ E_7&4032&15&30&6&1&3\\
  \typ E_8&13440&27&72&10&1&5.
  \end{array}
  \]}
  \begin{proof}\pf\ Exact integer evaluation gives \texttt{true} for
   \path{Computations.checkNoSolution} on $R\times R$, where a bad pair has dot
   product $-4$ and fails the specified row test. The lengths of the filtered
   ordered-pair lists are the three displayed counts. No real arithmetic is evaluated.~\qed\end{proof}
  \step{s3c4sound}{The row holds for every mathematical ordered star.}
  \begin{proof}\pf\ Apply \path{Computations.checkNoSolution_sound} to
   \stepref{s3c4census}, using the complete list identification in
   \stepref{s3c4lists}.~\qed\end{proof}
  \qedstep
  \begin{proof}\pf\ Adjacency is exactly a zero raw product; divide the length of each list $D(a,b)$
   by 2 to count the lines of $\Delta$. Thus \stepref{s3c4sound} gives the asserted parameters.~\qed\end{proof}
 \end{proof}
 \step{s3c4sigma}{\new\ The map $\sigma$ is an automorphism of $\typ D_\ell$, $\typ E_8$,
  $\typ E_7$ and $\typ E_6$, and $\sigma(\Za)=\Za'$, $\sigma(\Zb)=\Zb'$.}
 \begin{proof}
  \step{s3c4sigmaD}{$\sigma$ permutes $\Phi(\typ D_\ell)$ and the pair roots of model (ii).}
  \begin{proof}\pf\ A signed permutation of coordinates carries $s\Ze_i+t\Ze_j$ to a vector
   of the same form, in Definitions~\ref{def:3.1} and~\ref{def:3.2}.~\qed\end{proof}
  \step{s3c4sigmahalf}{$\sigma$ permutes the half roots of model (ii).}
  \begin{proof}\pf\ It permutes the entries of $\frac12\sum\epsilon(i)\Ze_i$ and changes
   two signs, so the product condition of Definition~\ref{def:3.2}(ii) is kept.~\qed\end{proof}
  \step{s3c4sigmacuts}{$\sigma$ permutes $\Phi(\typ E_7)$ and $\Phi(\typ E_6)$.}
  \begin{proof}\pf\ It fixes the coordinates 6, 7 and 8, so it keeps the equations of
   Definition~\ref{def:3.3}; use \stepref{s3c4sigmaD} and \stepref{s3c4sigmahalf}.~\qed\end{proof}
  \step{s3c4sigmastar}{$\sigma(\Za)=\Za'$ and $\sigma(\Zb)=\Zb'$.}
  \begin{proof}\pf\ $\sigma(\Ze_1-\Ze_2)=\Ze_2+\Ze_1$ and
   $\sigma(\Ze_2-\Ze_3)=-\Ze_1+\Ze_3$.~\qed\end{proof}
  \qedstep
  \begin{proof}\pf\ The map $\sigma$ is an orthogonal linear map; the four steps give the
   assertions.~\qed\end{proof}
 \end{proof}
 \step{s3c4labels}{For the ordered star $\Za'=\Ze_1+\Ze_2$,
  $\Zb'=-\Ze_1+\Ze_3$, the orthogonality graphs on $\Ga$ of $\typ E_6$ and $\typ E_7$
  are the lattice graph and the six-symbol disjointness graph.}
 \begin{proof}
  \pflet{$A_{\typ E_6}$ and $A_{\typ E_7}$ be the integer adjacency tables
   \texttt{e6Adjacency} and \texttt{e7Adjacency} of \path{Computations/C4E6Data.lean} and
   \path{Computations/C4E7Data.lean}, and let $A_{\rm lat}$ and $A_{\rm dis}$ be the
   tables of adjacency in one common coordinate on $\{0,1,2\}^2$ and of disjointness on
   $\binom{[6]}2$, in lexicographic order.}
  \step{s3c4Esixstar}{The roots $\Za',\Zb'$ are the marked star of the Lean term
   \leancite{LineSystem.E₆.hyp24}{\texttt{LineSystem.\allowbreak E\ensuremath{{}_6}.\allowbreak hyp24}}.}
  \begin{proof}\pf\ These are
   \leancite{LineSystem.E₆.hyp24_a}{\texttt{LineSystem.\allowbreak E\ensuremath{{}_6}.\allowbreak hyp24\_\allowbreak a}} and
   \leancite{LineSystem.E₆.hyp24_b}{\texttt{LineSystem.\allowbreak E\ensuremath{{}_6}.\allowbreak hyp24\_\allowbreak b}},
   in zero-based coordinates.~\qed\end{proof}
  \step{s3c4Esevenstar}{The roots $\Za',\Zb'$ are the marked star of the Lean term
   \leancite{LineSystem.E₇.hyp24}{\texttt{LineSystem.\allowbreak E\ensuremath{{}_7}.\allowbreak hyp24}}.}
  \begin{proof}\pf\ These are
   \leancite{LineSystem.E₇.hyp24_a}{\texttt{LineSystem.\allowbreak E\ensuremath{{}_7}.\allowbreak hyp24\_\allowbreak a}} and
   \leancite{LineSystem.E₇.hyp24_b}{\texttt{LineSystem.\allowbreak E\ensuremath{{}_7}.\allowbreak hyp24\_\allowbreak b}},
   in zero-based coordinates.~\qed\end{proof}
  \step{s3c4Esixtable}{The graph $\opk{Orth}(\Ga)$ of $\typ E_6$ for this star is
   isomorphic to the graph decoded from $A_{\typ E_6}$.}
  \begin{proof}\raggedright\pf\ This is
   \leancite{LineSystem.E₆.orthGraphIso}{\texttt{LineSystem.\allowbreak E\ensuremath{{}_6}.\allowbreak orthGraphIso}},
   by \stepref{s3c4Esixstar}. Its target
   \path{LineSystem.Hyp24.orthGraph} has adjacency at inner product 0, by
   \path{LineSystem.Hyp24.orthGraph_adj}.~\qed\end{proof}
  \step{s3c4Eseventable}{The graph $\opk{Orth}(\Ga)$ of $\typ E_7$ for this star is
   isomorphic to the graph decoded from $A_{\typ E_7}$.}
  \begin{proof}\raggedright\pf\ This is
   \leancite{LineSystem.E₇.orthGraphIso}{\texttt{LineSystem.\allowbreak E\ensuremath{{}_7}.\allowbreak orthGraphIso}},
   by \stepref{s3c4Esevenstar}, with the same adjacency
   \path{LineSystem.Hyp24.orthGraph_adj}.~\qed\end{proof}
  \step{s3c4grid}{The graph decoded from $A_{\typ E_6}$ is isomorphic to the lattice graph on $\{0,1,2\}^2$.}
  \begin{proof}\raggedright\pf\ Apply \path{Computations.nonempty_graphOfMatrix_iso} to
   \path{Computations.checkGraphIso} on $A_{\typ E_6}$, $A_{\rm lat}$ and a nine-entry permutation;
   in Lean it is the labelling \path{LineSystem.E₆.latticeLabel} of
   \path{Computations/C4Graphs.lean}.~\qed\end{proof}
  \step{s3c4duads}{The graph decoded from $A_{\typ E_7}$ is isomorphic to disjointness on $\binom{[6]}2$.}
  \begin{proof}\raggedright\pf\ Apply \path{Computations.nonempty_graphOfMatrix_iso} to
   \path{Computations.checkGraphIso} on $A_{\typ E_7}$, $A_{\rm dis}$ and a fifteen-entry permutation;
   in Lean it is the labelling \path{LineSystem.E₇.pairLabel} of the same file.~\qed\end{proof}
  \step{s3c4Esixlabel}{The $\typ E_6$ assertion holds.}
  \begin{proof}\pf\ Compose the isomorphisms of \stepref{s3c4Esixtable} and
   \stepref{s3c4grid}.~\qed\end{proof}
  \step{s3c4Esevenlabel}{The $\typ E_7$ assertion holds.}
  \begin{proof}\pf\ Compose the isomorphisms of \stepref{s3c4Eseventable} and
   \stepref{s3c4duads}.~\qed\end{proof}
  \qedstep
  \begin{proof}\pf\ Steps \stepref{s3c4Esixlabel} and \stepref{s3c4Esevenlabel} are the
   two assertions.~\qed\end{proof}
 \end{proof}
 \step{s3c4trans}{\new\ For $t\in\{6,7,8\}$ the isometries of $\re^8$ preserving
  $\Phi(\typ E_t)$ are transitive on the ordered stars of $\typ E_t$, as stated.}
 \begin{proof}\raggedright\pf\ This is
  \leancite{LineSystem.E₆.exists_isometry_star}{\texttt{LineSystem.\allowbreak E\ensuremath{{}_6}.\allowbreak exists\_\allowbreak isometry\_\allowbreak star}}
  and its analogues for $\typ E_7$ and $\typ E_8$, whose fixed star is $\Za',\Zb'$ by
  \leancite{LineSystem.E₈.hyp24_a}{\texttt{LineSystem.\allowbreak E\ensuremath{{}_8}.\allowbreak hyp24\_\allowbreak a}},
  \leancite{LineSystem.E₈.hyp24_b}{\texttt{LineSystem.\allowbreak E\ensuremath{{}_8}.\allowbreak hyp24\_\allowbreak b}}
  and their analogues; their trees of reflections are checked by
  \texttt{decide +kernel}.~\qed\end{proof}
 \step{s3c4ranks}{The ranks are $\ell$, $\ell$, 6, 7, 8, and $\typ A_1$ has one line and rank 1.}
 \begin{proof}\pf\ Read the ranks from Lemma~\ref{lem:sec3-spans}; the line of
  $\typ A_1$ is $\Ze_1-\Ze_2$ by Definition~\ref{def:3.1}.~\qed\end{proof}
 \qedstep
 \begin{proof}\pf\ Combine \stepref{s3c4closed}, \stepref{s3c4indec},
  \stepref{s3c4othermodels}, \stepref{s3c4Agraph}, \stepref{s3c4Dgraph},
  \stepref{s3c4exceptional}, \stepref{s3c4sigma}, \stepref{s3c4labels}, \stepref{s3c4trans},
  and \stepref{s3c4ranks}.
  None uses Theorem 2.13 or 3.5.~\qed\end{proof}
\end{proof}
\begingroup\raggedright
The two classical closure declarations are in \path{LineSystem/ClassicalModels.lean};
the three exceptional closure declarations are in \path{Computations/E8.lean}.
Computation C4 is proved in Lean, for the marked star and for every ordered star of the
exceptional models. The indecomposability theorems are in \path{Computations/C4.lean},
\path{Computations/C4Classical.lean} and \path{Computations/C4E6.lean} to
\path{Computations/C4E8.lean}; the marked stars \texttt{hyp24}, the isomorphisms
\texttt{orthGraphIso}, the counts \texttt{card\_\ensuremath{\Gamma}a} and \texttt{card\_\ensuremath{\Delta}\_div\_two}, and the
strongly regular parameters \texttt{orthGraph\_isSRGWith} are in
\path{Computations/C4An.lean}, \path{Computations/C4Dn.lean} and
\path{Computations/C4E6.lean} to \path{Computations/C4E8.lean}. The Lean star of $\typ A_\ell$
and of $\typ D_\ell$ is our $\Za,\Zb$ in zero-based coordinates, $\Ze_0-\Ze_1$, $\Ze_1-\Ze_2$.
The Lean star \texttt{hyp24} of $\typ E_6$, $\typ E_7$ and $\typ E_8$ is $\Za',\Zb'$, the image of
our star under $\sigma$; our star itself is \texttt{hyp24Simple}, and
\leancite{Computations.C4.marked_star_data}{\texttt{Computations.\allowbreak C4.\allowbreak marked\_\allowbreak star\_\allowbreak data}} records the table above for our star
in all five rows. For every ordered star, a reflection in a root of a star-closed system is an
automorphism, \leancite{LineSystem.rootReflection_mem_isometryGroup}{\texttt{LineSystem.\allowbreak rootReflection\_\allowbreak mem\_\allowbreak isometryGroup}} in
\path{LineSystem/Reflections.lean}, and two reflection trees for each exceptional model, checked
by \texttt{decide +kernel} in \path{Computations/C4E6Stars.lean} to
\path{Computations/C4E8Stars.lean}, give transitivity on ordered stars:
\leancite{LineSystem.E₆.exists_isometry_star}{\texttt{LineSystem.\allowbreak E\ensuremath{{}_6}.\allowbreak exists\_\allowbreak isometry\_\allowbreak star}}, \leancite{LineSystem.E₇.exists_isometry_star}{\texttt{LineSystem.\allowbreak E\ensuremath{{}_7}.\allowbreak exists\_\allowbreak isometry\_\allowbreak star}} and
\leancite{LineSystem.E₈.exists_isometry_star}{\texttt{LineSystem.\allowbreak E\ensuremath{{}_8}.\allowbreak exists\_\allowbreak isometry\_\allowbreak star}}. For every \texttt{H : Hyp24}
with \texttt{H.S} one of the three models, \leancite{LineSystem.Hyp24.orthGraph_isSRGWith_of_S_eq_E₆}{\texttt{LineSystem.\allowbreak Hyp24.\allowbreak orthGraph\_\allowbreak isSRGWith\_\allowbreak of\_\allowbreak S\_\allowbreak eq\_\allowbreak E\ensuremath{{}_6}}},
\leancite{LineSystem.Hyp24.card_Γa_of_S_eq_E₆}{\texttt{LineSystem.\allowbreak Hyp24.\allowbreak card\_\allowbreak \ensuremath{\Gamma}a\_\allowbreak of\_\allowbreak S\_\allowbreak eq\_\allowbreak E\ensuremath{{}_6}}} and
\leancite{LineSystem.Hyp24.card_Δ_of_S_eq_E₆}{\texttt{LineSystem.\allowbreak Hyp24.\allowbreak card\_\allowbreak \ensuremath{\Delta}\_\allowbreak of\_\allowbreak S\_\allowbreak eq\_\allowbreak E\ensuremath{{}_6}}}, with their $\typ E_7$ and $\typ E_8$
analogues, give the exceptional rows; Lean counts $\Delta$ as $12,30,72$ roots, twice the
entries of the last column. The bundle \leancite{Computations.C4.exceptional_star_data}{\texttt{Computations.\allowbreak C4.\allowbreak exceptional\_\allowbreak star\_\allowbreak data}}
collects them. The graph labels are the isomorphisms
\leancite{LineSystem.E₆.latticeGraphIso}{\texttt{LineSystem.\allowbreak E\ensuremath{{}_6}.\allowbreak latticeGraphIso}}, from the line graph of $K_{3,3}$, and
\leancite{LineSystem.E₇.triangularGraphComplIso}{\texttt{LineSystem.\allowbreak E\ensuremath{{}_7}.\allowbreak triangularGraphComplIso}}, from the complement of the line
graph of $K_6$, both for the star $\Za',\Zb'$, in \path{Computations/C4Graphs.lean}; their every-star
forms are \leancite{LineSystem.Hyp24.nonempty_latticeGraph_iso_of_S_eq_E₆}{\texttt{LineSystem.\allowbreak Hyp24.\allowbreak nonempty\_\allowbreak latticeGraph\_\allowbreak iso\_\allowbreak of\_\allowbreak S\_\allowbreak eq\_\allowbreak E\ensuremath{{}_6}}} and
\leancite{LineSystem.Hyp24.nonempty_triangularGraph_compl_iso_of_S_eq_E₇}{\texttt{LineSystem.\allowbreak Hyp24.\allowbreak nonempty\_\allowbreak triangularGraph\_\allowbreak compl\_\allowbreak iso\_\allowbreak of\_\allowbreak S\_\allowbreak eq\_\allowbreak E\ensuremath{{}_7}}}. In Lean the
transitivity replaces the transfer by $\sigma$. All these isomorphisms are for the orthogonality
graph; the bridge \path{LineSystem.IntegerModel.adjMatrix_eq_dotAdjacency} builds the graph at
inner product 1, and we do not use it. The script \path{blueprint/checks/sec3/check_sec3.py}
recomputes the finite data; no Lean proof evaluates a spectrum.\par
\endgroup
The source lists have the membership theorems
\leancite{LineSystem.E₆.mem_rootList}{\texttt{LineSystem.\allowbreak E\ensuremath{{}_6}.\allowbreak mem\_\allowbreak rootList}}, \leancite{LineSystem.E₇.mem_rootList}{\texttt{LineSystem.\allowbreak E\ensuremath{{}_7}.\allowbreak mem\_\allowbreak rootList}}, and
\leancite{LineSystem.E₈viaD₈.mem_rootList}{\texttt{LineSystem.\allowbreak E\ensuremath{{}_8}viaD\ensuremath{{}_8}.\allowbreak mem\_\allowbreak rootList}}. Sorting does not change membership;
the lists $R,L,D$ in C4 contain each raw vector once.

\begin{namedresult}{Lemma}{(independence of the perpendicular choices)}\label{lem:sec3-choices}
 \lean{LineSystem.E₈.nonempty_iso_E₇_orthogonalTo}
 \lean{LineSystem.E₈.nonempty_iso_E₆_orthogonalTo}
 \lean{LineSystem.E₈.exists_isometry_of_ordered_star}
 \lean{LineSystem.E₈.exists_isometry_of_root}
 \lean{LineSystem.E₈viaA₈.exists_isometry_of_ordered_stars}
 \lean{LineSystem.E₈viaA₈.exists_isometry_of_roots}
 \lean{LineSystem.E₈viaA₈.nonempty_iso_E₇_orthogonalTo}
 \lean{LineSystem.E₈viaA₈.nonempty_iso_E₆_orthogonalTo}
 \lean{LineSystem.E₈viaSteiner.exists_isometry_of_ordered_stars}
 \lean{LineSystem.E₈viaSteiner.exists_isometry_of_ordered_star}
 \lean{LineSystem.E₈viaSteiner.exists_isometry_of_roots}
 \lean{LineSystem.E₈viaSteiner.exists_isometry_of_root}
 \lean{LineSystem.E₈viaSteiner.nonempty_iso_E₇_orthogonalTo}
 \lean{LineSystem.E₈viaSteiner.nonempty_iso_E₆_orthogonalTo}
 \lean{LineSystem.T₇.nonempty_iso_E₇}\lean{LineSystem.T₆.nonempty_iso_E₆}
 \uses{def:3.2,def:3.3,comp:C2,comp:C3,comp:C4,lem:2.10,lem:2.12,def:indecomposable,facts:inner}
 \assume{$\typ E_8$ is the default model in $\re^8$; $\Zr\in\Phi(\typ E_8)$ and,
 for the star assertion, $\Zs\in\Phi(\typ E_8)$ with $(\Zr,\Zs)=-1$.}
 \prove{There is an automorphism of $\typ E_8$, induced by a linear isometry of
 $\re^8$, carrying $(\Zr,\Zs)$ to $(\Zr_0,\Zs_0)$.
 Every root $\Zr$ also admits such an automorphism carrying it to $\Zr_0$,
 without an initially chosen partner. Thus the cuts perpendicular to any line
 or any star are isometric to $\typ E_7$ or $\typ E_6$, respectively.
 The same assertions hold in the other two models, and $T_7\cong \typ E_7$, $T_6\cong \typ E_6$.}
\end{namedresult}
\begin{proof}
 \pfsketch\ \repair\ We fill Gap 5 before classification. C4 and Lemma 2.10
 identify the marked graphs; Lemma 2.12 extends their isomorphism to the ordered stars.
 \step{s3choicepair}{Any two ordered stars of $\typ E_8$ are carried to each other by an automorphism.}
 \begin{proof}
  \step{s3choicegraphs}{Their graphs on $\Ga$ are isomorphic.}
  \begin{proof}\pf\ C4 gives order 27, degree 10, and star-closure and
   indecomposability for each instance. Apply Lemma~\ref{lem:2.10}; the positive
   degree excludes an empty graph and degree $10<26$ excludes a universal vertex.~\qed\end{proof}
  \qedstep
  \begin{proof}\pf\ Lemma~\ref{lem:2.12} extends this graph isomorphism to a
   surjective span isometry carrying the ordered stars to each other.~\qed\end{proof}
 \end{proof}
 \step{s3choiceroot}{Every root has a partner of inner product $-1$.}
 \begin{proof}
  \step{s3choicenonorth}{Some root outside $\{\Zr,-\Zr\}$ has nonzero product with $\Zr$.}
  \begin{proof}\pf\ Otherwise this antipodal pair and its nonempty complement
   would be an orthogonal partition. This contradicts C4; the complement is
   nonempty since C2 gives 120 lines.~\qed\end{proof}
  \qedstep
  \begin{proof}\pf\ The product in \stepref{s3choicenonorth} is $1$ or $-1$
   by Section~\ref{facts:inner}; negate the partner when the product is 1.~\qed\end{proof}
 \end{proof}
 \step{s3choicecuts}{The default-model perpendicular cuts have the asserted isomorphisms.}
 \begin{proof}
  \step{s3choiceSeven}{The cut perpendicular to $\Zr$ maps onto $\typ E_7$.}
  \begin{proof}\pf\ Use \stepref{s3choiceroot} and \stepref{s3choicepair} to send
   $\Zr$ to $\Zr_0$; an isometry preserves the equation $(\Zx,\Zr)=0$.~\qed\end{proof}
  \step{s3choiceSix}{The cut perpendicular to the ordered star maps onto $\typ E_6$.}
  \begin{proof}\pf\ Step \stepref{s3choicepair} preserves the two equations of
   orthogonality to $\Zr,\Zs$; orthogonality to $-\Zr-\Zs$ is their sum.~\qed\end{proof}
  \qedstep
  \begin{proof}\pf\ These are the two cuts in Definition~\ref{def:3.3}.~\qed\end{proof}
 \end{proof}
 \step{s3choiceSteiner}{In model (iii), $T_7$ and $T_6$ are cuts of the required kinds.}
 \lean{LineSystem.T₇_eq_orthogonalTo}\lean{LineSystem.T₆_eq_orthogonalTo}
 \lean{LineSystem.inner_t₇Normal_t₆Normal}
 \begin{proof}\pf\ Take $\Za=\sqrt2\Ze_8$ and
  $\Zb=-(\Ze_5+\Ze_6+\Ze_7+\Ze_8)/\sqrt2$.
  Here $\Za$ is an axis root and $\Zb$ is a root by block 5678, and their product is
  $-1$. Their two perpendicular
  equations are exactly those defining $T_7,T_6$.~\qed\end{proof}
 \qedstep
 \begin{proof}\pf\ C3 carries any root or ordered star in either other model to
  the default model. Apply \stepref{s3choicecuts}, including the cuts identified
  by \stepref{s3choiceSteiner}.~\qed\end{proof}
\end{proof}

\begin{cgss}{Lemma}{3.4}\label{lem:3.4}
 \lean{LineSystem.numLines_An}\lean{LineSystem.numLines_Dn}
 \lean{LineSystem.numLines_E₈}\lean{LineSystem.numLines_E₈viaA₈}
 \lean{LineSystem.numLines_E₈viaSteiner}\lean{LineSystem.numLines_E₇}\lean{LineSystem.numLines_E₆}
 \uses{def:3.1,def:3.2,def:3.3,comp:C2,lem:sec3-choices}
 \assume{The model parameters are $\ell\geq1$ for $\typ A_\ell$ and $\ell\geq4$ for
 $\typ D_\ell$; the exceptional systems have the definitions above.}
 \prove{\repair\ With $|\typ A_\ell|$ in place of the printed $|A|$ (Error 13),
 \[
 |\typ A_\ell|=\tfrac12\ell(\ell+1),\quad |\typ D_\ell|=\ell(\ell-1),\quad
 |\typ E_8|=120,\quad |\typ E_7|=63,\quad |\typ E_6|=36.
 \]
 These are counts of lines, independent of the model and of the perpendicular choices.}
\end{cgss}
\begin{proof}
 \pfsketch\ We will use the signed counts of C2, and then justify the paper's
 use of model (iii) for the perpendicular systems.
 \step{s334classical}{The first two formulas hold.}
 \begin{proof}\pf\ Apply the first two counts of C2 in the stated parameter ranges.~\qed\end{proof}
 \step{s334eight}{All three definitions of $\typ E_8$ have 120 lines.}
 \begin{proof}\pf\ C2 gives $36+84=56+64=8+8\cdot14=120$.~\qed\end{proof}
 \step{s334seven}{Every perpendicular-line system has 63 lines.}
 \begin{proof}\pf\ Lemma~\ref{lem:sec3-choices} identifies it with $T_7$;
  C2 gives $7+8\cdot7=63$.~\qed\end{proof}
 \step{s334six}{Every perpendicular-star system has 36 lines.}
 \begin{proof}\pf\ Lemma~\ref{lem:sec3-choices} identifies it with $T_6$;
  C2 gives $4+8+4\cdot6=36$.~\qed\end{proof}
 \qedstep
 \begin{proof}\pf\ These four steps cover all five types and all allowed choices.
  The seven named count declarations already exist in \path{LineSystem/Models.lean}.~\qed\end{proof}
\end{proof}

\begin{cgss}{Theorem}{3.5}\label{thm:3.5}
 \lean{LineSystem.classification}\lean{LineSystem.RootType}
 \lean{LineSystem.RootType.model}\lean{LineSystem.RootType.rank}
 \uses{thm:2.13,def:line-system,def:3.1,def:3.2,def:3.3,comp:C2,comp:C3,comp:C4,lem:3.4,lem:sec3-spans,lem:sec3-transport,lem:sec3-restrict}
 \assume{\begin{pfenum}
  \item $V$ is a finite-dimensional real inner product space.
  \item $S$ is a nonempty line system in $V$, and $W=\operatorname{span}\Phi(S)$.
 \end{pfenum}}
 \prove{$S$ is indecomposable and star-closed if and only if it is isomorphic to
 exactly one of
 $\typ A_\ell$ ($\ell\geq1$), $\typ D_\ell$ ($\ell\geq4$), $\typ E_6$, $\typ E_7$, $\typ E_8$.
 The isomorphism is induced by a surjective linear isometry of spans.
 All three definitions of $\typ E_8$ give that same type.}
 \new\ In Lean the list is an index type \texttt{LineSystem.RootType}, with
 constructors for $\typ A_\ell$ with $1\leq\ell$ and $\typ D_\ell$ with $4\leq\ell$, and
 three constants for $\typ E_6$, $\typ E_7$, $\typ E_8$. Each type $t$ has an ambient
 dimension \texttt{t.dim}, a model \texttt{t.model}, a line system in the Euclidean space
 of that dimension, and a rank \texttt{t.rank}. The Lean statement then says that $S$ is
 indecomposable and star-closed if and only if there is a unique $t$ with an isomorphism
 from $S$ to \texttt{t.model}. Theorem 3.6 uses the same type, with $\typ A_7$, $\typ A_8$
 and $\typ D_8$ excluded.
\end{cgss}
\begin{proof}
 \pfsketch\ The models realize the five possibilities of Theorem 2.13. Its inputs
 from this section are C2, C4, and the span calculations, which have all been
 established without either classification theorem. We will include the one-line case.
 \new\ The uniqueness assertion of Theorem 2.13 is stated for systems spanning
 $\re^\ell$. We will restrict $S$ and the realizing model to their spans, map both spans
 to $\re^\ell$ through orthonormal bases, apply uniqueness there, and compose.
 \step{s335forward}{\assume{$S$ is indecomposable and star-closed.}
  \prove{There are a listed model $Y$ and a surjective linear isometry
  $F:W\to\operatorname{span}\Phi(Y)$ with $F(\Phi(S))=\Phi(Y)$.}}
 \begin{proof}
  \pflet{$\ell=\dim W$, and let $S|_W$ be the restriction of $S$ to $W$ in
   Lemma~\ref{lem:sec3-restrict}.}
  \step{s335restrict}{$S|_W$ satisfies the hypotheses of Theorem~\ref{thm:2.13},
   with $W$ as the ambient space of dimension $\ell$.}
  \begin{proof}\pf\ Lemma~\ref{lem:sec3-restrict} makes $S|_W$ a line system spanning
   $W$, nonempty by \asm[2]{0} and indecomposable and star-closed by \asm{1}; $W$ is
   finite-dimensional by \asm[1]{0}.~\qed\end{proof}
  \step{s335possibilities}{$S|_W$ lies in one of the five cases of Theorem 2.13.}
  \begin{proof}\pf\ Apply Theorem~\ref{thm:2.13} to \stepref{s335restrict}.~\qed\end{proof}
  \step{s335realized}{A listed model $Y$ of rank $\ell$ has the number of lines of that
   case and, when that case has more than one line, an ordered star with the data on
   $\Ga$ and $\Delta$ of that case.}
  \begin{proof}
   \step{s335A}{\case{Case 1 of Theorem 2.13.}}
   \begin{proof}\pf\ Take $Y=\typ A_\ell$: C4 gives rank $\ell$, the empty graph on
    $\ell-2$ vertices for $\ell\geq2$, and one line for $\ell=1$; Lemma~\ref{lem:3.4}
    gives $\ell(\ell+1)/2$ lines.~\qed\end{proof}
   \step{s335D}{\case{Case 2 of Theorem 2.13, so $\ell\geq4$.}}
   \begin{proof}\pf\ Take $Y=\typ D_\ell$: C4 gives rank $\ell$ and the friendship graph
    of $\ell-3$ triangles; Lemma~\ref{lem:3.4} gives $\ell(\ell-1)$ lines.~\qed\end{proof}
   \step{s335Eeight}{\case{Case 3 of Theorem 2.13, so $\ell=8$.}}
   \begin{proof}\pf\ Take $Y=\typ E_8$: C4 gives rank 8, $|\Ga|=27$ and
    $|\Delta|/2=36$; Lemma~\ref{lem:3.4} gives 120 lines.~\qed\end{proof}
   \step{s335Eseven}{\case{Case 4 of Theorem 2.13, so $\ell=7$.}}
   \begin{proof}\pf\ Take $Y=\typ E_7$: C4 gives rank 7, $|\Ga|=15$ and
    $|\Delta|/2=15$; Lemma~\ref{lem:3.4} gives 63 lines.~\qed\end{proof}
   \step{s335Esix}{\case{Case 5 of Theorem 2.13, so $\ell=6$.}}
   \begin{proof}\pf\ Take $Y=\typ E_6$: C4 gives rank 6, $|\Ga|=9$ and
    $|\Delta|/2=6$; Lemma~\ref{lem:3.4} gives 36 lines.~\qed\end{proof}
   \qedstep
   \begin{proof}\pf\ These are the five cases in \stepref{s335possibilities}.~\qed\end{proof}
  \end{proof}
  \pflet{$W_Y=\operatorname{span}\Phi(Y)$, and let $Y|_{W_Y}$ be the restriction of $Y$
   to $W_Y$ in Lemma~\ref{lem:sec3-restrict}.}
  \step{s335modelhyp}{$Y|_{W_Y}$ satisfies the hypotheses of
   Theorem~\ref{thm:2.13}, with $W_Y$ as the ambient space of dimension $\ell$.}
  \begin{proof}
   \step{s335modelprops}{$Y$ is nonempty, indecomposable and star-closed.}
   \begin{proof}\pf\ This is the first assertion of C4.~\qed\end{proof}
   \step{s335modeldim}{$\dim W_Y=\ell$.}
   \begin{proof}\pf\ The rank of $Y$ is $\ell$ by \stepref{s335realized}.~\qed\end{proof}
   \qedstep
   \begin{proof}\pf\ Lemma~\ref{lem:sec3-restrict} makes $Y|_{W_Y}$ a line system
    spanning $W_Y$ with the properties of \stepref{s335modelprops}; its dimension is
    \stepref{s335modeldim}.~\qed\end{proof}
  \end{proof}
  \step{s335samecase}{$Y|_{W_Y}$ lies in the same case of Theorem 2.13 as $S|_W$.}
  \begin{proof}\pf\ Lemma~\ref{lem:sec3-restrict} carries the number of lines and the
   marked-star data of \stepref{s335realized} from $Y$ to $Y|_{W_Y}$, whose dimension is
   $\ell$ by \stepref{s335modelhyp}; these are the data of the case of
   \stepref{s335possibilities}.~\qed\end{proof}
  \step{s335coords}{There are surjective linear isometries $\iota:W\to\re^\ell$
   and $\iota_Y:W_Y\to\re^\ell$; the sets $\iota(\Phi(S))$ and $\iota_Y(\Phi(Y))$ are the
   root sets of line systems $S'$ and $Y'$ in $\re^\ell$, which satisfy the hypotheses of
   Theorem~\ref{thm:2.13} and lie in the case of \stepref{s335possibilities}.}
  \begin{proof}
   \step{s335basis}{The isometries $\iota$ and $\iota_Y$ exist.}
   \begin{proof}\pf\ Take the representation \path{OrthonormalBasis.repr} of
    \path{stdOrthonormalBasis} of $W$ and of $W_Y$, reindexed to $\operatorname{Fin}\ell$
    by \path{OrthonormalBasis.reindex} along \path{finCongr}, using $\dim W=\ell$ and
    \stepref{s335modelhyp}.~\qed\end{proof}
   \step{s335imagelines}{$S'$ and $Y'$ are line systems in $\re^\ell$, and $\iota$ and
    $\iota_Y$ restrict to bijections of root sets that preserve inner products.}
   \begin{proof}\pf\ A linear isometry is injective, commutes with negation and
    preserves inner products, by \path{LinearIsometry.inner_map_map}; the conditions of
    the line-system definition in Section~\ref{def:line-system} are of these kinds.~\qed\end{proof}
   \step{s335imagespan}{$\Phi(S')$ and $\Phi(Y')$ span $\re^\ell$.}
   \begin{proof}\pf\ By \path{Submodule.span_image}, the span of an image is the image
    of the span, which is the whole of $\re^\ell$ for a surjective map; use
    \stepref{s335restrict} and \stepref{s335modelhyp}.~\qed\end{proof}
   \qedstep
   \begin{proof}\pf\ Lemma~\ref{lem:sec3-transport}, applied to the bijections of
    \stepref{s335imagelines}, carries nonemptiness, indecomposability, star-closure, the
    number of lines and the marked-star data from $S|_W$ and $Y|_{W_Y}$, which have them
    by \stepref{s335restrict}, \stepref{s335possibilities}, \stepref{s335modelhyp} and
    \stepref{s335samecase}; \stepref{s335imagespan} gives spanning.~\qed\end{proof}
  \end{proof}
  \step{s335orthogonal}{There is an orthogonal transformation $O$ of $\re^\ell$ with
   $O(\Phi(S'))=\Phi(Y')$.}
  \begin{proof}\pf\ Apply the uniqueness assertion of Theorem~\ref{thm:2.13}, for the
   ambient space $\re^\ell$, to \stepref{s335coords}.~\qed\end{proof}
  \qedstep
  \begin{proof}\pf\ Put $F=\iota_Y^{-1}\circ O\circ\iota$, a surjective linear isometry
   $W\to W_Y$ by \path{LinearIsometryEquiv.symm} and \path{LinearIsometryEquiv.trans};
   \stepref{s335coords} and \stepref{s335orthogonal} give
   $F(\Phi(S))=\Phi(Y)$.~\qed\end{proof}
 \end{proof}
 \step{s335reverse}{\assume{$S$ is isomorphic to a listed model.}
  \prove{$S$ is indecomposable and star-closed.}}
 \begin{proof}\pf\ C4 gives these properties for the model; apply
  Lemma~\ref{lem:sec3-transport}.~\qed\end{proof}
 \step{s335unique}{Different listed types are not isomorphic.}
 \begin{proof}
  \step{s335rank}{Isomorphic models have the same rank.}
  \begin{proof}\pf\ Apply Lemma~\ref{lem:sec3-transport}.~\qed\end{proof}
  \step{s335AD}{At a common rank $\ell\geq4$, the $\typ D$ count exceeds the $\typ A$ count by $\ell(\ell-3)/2>0$.}
  \begin{proof}\pf\ Subtract the formulas in Lemma~\ref{lem:3.4}.~\qed\end{proof}
  \step{s335exceptionalcounts}{At ranks $6,7,8$ the respective $(\typ A,\typ D,\typ E)$ counts are
   $(21,30,36)$, $(28,42,63)$, $(36,56,120)$, with strict inequalities in each row.}
  \begin{proof}\pf\ Substitute those ranks in Lemma~\ref{lem:3.4}.~\qed\end{proof}
  \qedstep
  \begin{proof}\pf\ Line count is an isomorphism invariant by
   Lemma~\ref{lem:sec3-transport}; \stepref{s335rank}--\stepref{s335exceptionalcounts}
   distinguish every pair of listed types.~\qed\end{proof}
 \end{proof}
 \step{s335atmost}{The system $S$ cannot be isomorphic to two different listed types.}
 \begin{proof}\pf\ Compose one span isometry with the inverse of the other from
  Lemma~\ref{lem:sec3-transport}. The resulting isometry between the two models
  contradicts \stepref{s335unique}.~\qed\end{proof}
 \qedstep
 \begin{proof}\pf\ Steps \stepref{s335forward}, \stepref{s335reverse}, and
  \stepref{s335atmost} prove the equivalence and uniqueness. C3 identifies the
  three $\typ E_8$ models independently.~\qed\end{proof}
\end{proof}

\begin{namedresult}{Lemma}{(a maximal extension in the span)}\label{lem:sec3-extension}
 \lean{LineSystem.exists_maximal_extension_span}
 \uses{def:line-system,def:indecomposable,def:2.2,lem:2.3,facts:inner}
 \assume{$S$ is a nonempty indecomposable line system in a finite-dimensional
 real inner product space, $W=\operatorname{span}\Phi(S)$, and $S$ is not maximal in $W$.}
 \prove{There is a line system $T$ in $W$ with
 $\Phi(S)\subsetneq\Phi(T)$, spanning $W$, that is maximal in $W$,
 indecomposable, and star-closed.}
\end{namedresult}
\begin{proof}
 \pfsketch\ \repair\ For the second omission in Gap 13 we will bound all candidate
 roots by their finite signatures against $\Phi(S)$, maximize the size of an extension,
 and prove indecomposability using the equality of spans.
 \pflet{$K=\{\Zx\in W:(\Zx,\Zx)=2,\ (\Zx,\Zr)\in\{-2,-1,0,1,2\}
 \text{ for every }\Zr\in\Phi(S)\}$.}
 \step{s3extfinite}{$K$ is finite.}
 \begin{proof}
  \step{s3extsignature}{The signature map $\Zx\mapsto((\Zx,\Zr))_{\Zr\in\Phi(S)}$ is injective on $K$.}
  \begin{proof}
   \step{s3extdifference}{Equal signatures make $\Zx-\Zy$ orthogonal to $W$.}
   \begin{proof}\pf\ Use bilinearity and the finite-span lemma in Section~\ref{facts:inner}.~\qed\end{proof}
   \qedstep
   \begin{proof}\pf\ Since $\Zx-\Zy\in W$, \stepref{s3extdifference} gives zero
    squared norm, so $\Zx=\Zy$ by Section~\ref{facts:inner}.~\qed\end{proof}
  \end{proof}
  \qedstep
  \begin{proof}\pf\ The signatures lie in a set of $5^{|\Phi(S)|}$ functions.
   Apply \path{Set.Finite.of_injOn} to \stepref{s3extsignature}.~\qed\end{proof}
 \end{proof}
 \step{s3extchoose}{Choose a largest-cardinality subset $\Psi\subseteq K$ containing
  $\Phi(S)$ that is the signed root set of a line system.}
 \begin{proof}\pf\ The eligible subsets form a nonempty finite family: it contains
  $\Phi(S)$ and is contained in the powerset of the finite $K$. Apply
  \path{Finset.exists_max_image} to their cardinalities.~\qed\end{proof}
 \pflet{$T$ be this line system with root set $\Psi$.}
 \step{s3extmaximal}{$T$ is maximal within $W$.}
 \begin{proof}
  \step{s3extcandidate}{\assume{$\Zx\in W$ has squared norm 2 and allowed products with $\Psi$.}
   \prove{$\Zx\in\Psi$.}}
  \begin{proof}
   \step{s3extinK}{$\Zx,-\Zx\in K$.}
   \begin{proof}\pf\ The candidate products include those with $\Phi(S)\subseteq\Psi$;
    the allowed product set is closed under negation.~\qed\end{proof}
   \step{s3extadjoin}{$\Psi\cup\{\Zx,-\Zx\}$ is eligible in \stepref{s3extchoose}.}
   \begin{proof}\pf\ Old-new products are allowed by \asm{2}, and the new-new
    products are $\pm2$; use \stepref{s3extinK} for containment in $K$.~\qed\end{proof}
   \qedstep
   \begin{proof}\pf\ Maximal cardinality in \stepref{s3extchoose} forbids a proper
    enlargement by \stepref{s3extadjoin}.~\qed\end{proof}
  \end{proof}
  \qedstep
  \begin{proof}\pf\ Step \stepref{s3extcandidate} is Definition~\ref{def:2.2}.~\qed\end{proof}
 \end{proof}
 \step{s3extspan}{$\operatorname{span}\Psi=W$.}
 \begin{proof}\pf\ $\Phi(S)\subseteq\Psi\subseteq W$ and $\Phi(S)$ spans $W$.~\qed\end{proof}
 \step{s3extindec}{$T$ is indecomposable.}
 \begin{proof}
  \step{s3extpartition}{\assume{$P\cup Q=\Psi$ and $P\perp Q$.}
   \prove{$P=\varnothing$ or $Q=\varnothing$.}}
  \begin{proof}
   \step{s3extoldpart}{All roots of $S$ lie in one of $P,Q$.}
   \begin{proof}\pf\ Intersect the partition with $\Phi(S)$ and use its indecomposability.~\qed\end{proof}
   \step{s3extempty}{The other part is empty.}
   \begin{proof}\pf\ Each root there is orthogonal to the spanning $\Phi(S)$,
    thus to itself by Section~\ref{facts:inner}; its norm would be zero.~\qed\end{proof}
   \qedstep
   \begin{proof}\pf\ Use \stepref{s3extempty}.~\qed\end{proof}
  \end{proof}
  \qedstep
  \begin{proof}\pf\ Apply the indecomposability definition in Section~\ref{def:indecomposable} to \stepref{s3extpartition}.~\qed\end{proof}
 \end{proof}
 \step{s3extproper}{$\Phi(S)\subsetneq\Psi$.}
 \begin{proof}\pf\ Equality would make $S$ maximal by \stepref{s3extmaximal},
  contrary to the hypothesis.~\qed\end{proof}
 \step{s3extclosed}{$T$ is star-closed.}
 \begin{proof}\pf\ Lemma~\ref{lem:2.3} applies to \stepref{s3extmaximal}.~\qed\end{proof}
 \qedstep
 \begin{proof}\pf\ Combine \stepref{s3extspan}, \stepref{s3extindec},
  \stepref{s3extproper}, and \stepref{s3extclosed}.~\qed\end{proof}
\end{proof}
The finiteness and maximum declarations used here are in
\path{Mathlib/Data/Set/Finite/Basic.lean} and \path{Mathlib/Data/Finset/Max.lean}.

\begin{namedresult}{Lemma}{(marked inclusions)}\label{lem:sec3-marked-inclusion}
 \lean{LineSystem.Hyp24.orthGraph_embedding_of_subset}
 \uses{hyp:2.4,lem:sec3-transport}
 \assume{$S,T$ are indecomposable star-closed line systems and a linear isometry
 of spans sends $\Phi(S)$ into $\Phi(T)$. Choose $\Za,\Zb\in\Phi(S)$ with
 $(\Za,\Zb)=-1$, and mark $T$ by their images.}
 \prove{The isometry induces an injection from the source $\Ga$ to the target
 $\Ga$ that preserves both adjacency and nonadjacency in their orthogonality graphs.}
\end{namedresult}
\begin{proof}
 \pfsketch\ \new\ We state the induced-embedding direction needed for C5 separately
 from the isomorphism assertion of Lemma 2.12. We will retain the same marked star
 and use the defining inner products.
 \step{s3markmembership}{A source member of $\Ga$ maps into the target $\Ga$.}
 \begin{proof}\pf\ The three defining products of Hypothesis~\ref{hyp:2.4}
  are preserved, and the image is a target root.~\qed\end{proof}
 \step{s3markinjection}{The induced map is injective.}
 \begin{proof}\pf\ A linear isometry is injective.~\qed\end{proof}
 \step{s3markedges}{For two distinct source vertices, the product is zero if and only if the image product is zero.}
 \begin{proof}\pf\ The isometry preserves that product.~\qed\end{proof}
 \qedstep
 \begin{proof}\pf\ The three steps give an induced graph embedding.~\qed\end{proof}
\end{proof}

\begin{namedresult}{Computation}{C5}\label{comp:C5}
 \lean{Computations.C5.inclusions_and_obstructions}
 \lean{LineSystem.not_embeds_An_Dn}
 \uses{def:3.1,def:3.2,def:3.3,comp:C3,comp:C4,lem:3.4,lem:sec3-spans,lem:sec3-choices,lem:sec3-marked-inclusion,lem:2.10}
 \assume{All inclusions are isometric inclusions of the signed models, in their spans.}
 \prove{There are proper inclusions $\typ A_8\subset \typ E_8$, $\typ D_8\subset \typ E_8$,
 and $\typ A_7\subset \typ E_7$, all of equal rank. There are no inclusions
 $\typ A_\ell\subset \typ D_\ell$ for $\ell\geq4$, $\typ A_6\subset \typ E_6$,
 $\typ D_6\subset \typ E_6$, or $\typ D_7\subset \typ E_7$.}
\end{namedresult}
\begin{proof}
 \pfsketch\ We will exhibit the inclusions, then obstruct induced marked graphs
 using coclique sizes and degrees. The $\typ A_\ell$ obstruction fills the first omission in Gap 13.
 \step{s3c5uniform}{The target marked graphs have the asserted types for every ordered star.}
 \begin{proof}
  \step{s3c5Dtransitive}{Signed coordinate permutations of $\typ D_\ell$ are transitive on ordered root pairs of product $-1$.}
  \begin{proof}
   \step{s3c5Dsupport}{Such roots share exactly one support position, with opposite signs there.}
   \begin{proof}\pf\ Equal supports give product $-2,0,2$; disjoint supports give
    0; a one-position intersection contributes its single sign product.~\qed\end{proof}
   \qedstep
   \begin{proof}\pf\ Send the unshared position of the first root to 1, the shared
    position to 2, and the unshared position of the second to 3. Choose signs
    to obtain $\Ze_1-\Ze_2,\Ze_2-\Ze_3$; extend the coordinate permutation and signs arbitrarily.~\qed\end{proof}
  \end{proof}
  \step{s3c5Euniform}{Every $\typ E_6$ marked graph is the lattice graph; every $\typ E_7$ marked graph is the disjointness graph on six symbols.}
  \begin{proof}\pf\ C4 gives constant order and parameters for every star.
   Lemma~\ref{lem:2.10} identifies its graph with the fixed-star graph, whose
   labelled type is certified in C4.~\qed\end{proof}
  \qedstep
  \begin{proof}\pf\ For $\typ D_\ell$ apply C4 after \stepref{s3c5Dtransitive}; for
   the two exceptional targets use \stepref{s3c5Euniform}.~\qed\end{proof}
 \end{proof}
 \step{s3c5Aeight}{$\typ A_8$ embeds properly in $\typ E_8$ with equal rank.}
 \begin{proof}\pf\ It is a subset of model (i); apply C3. Lemma~\ref{lem:sec3-spans}
  gives rank 8 on both sides, and Lemma~\ref{lem:3.4} gives $36<120$ lines.~\qed\end{proof}
 \step{s3c5Deight}{$\typ D_8$ embeds properly in $\typ E_8$ with equal rank.}
 \begin{proof}\pf\ Definition~\ref{def:3.2}(ii) includes these roots.
  Their ranks are 8 and their counts satisfy $56<120$ by the same two lemmas.~\qed\end{proof}
 \step{s3c5Aseven}{\new\ $\typ A_7$ embeds properly in $\typ E_7$ with equal rank. An explicit
  coordinate inclusion replaces the paper's embedding of marked graphs.}
 \begin{proof}
  \pflet{$\Zh=\frac12\mathbf1\in\re^8$.}
  \step{s3c5Ah}{$\typ A_7\subset \typ E_8\cap\Zh^\perp$.}
  \begin{proof}\pf\ Its roots are $\Ze_i-\Ze_j$ in $\re^8$, which are $\typ D_8$
   roots with product zero with $\Zh$. The all-positive half root $\Zh$ belongs to $\typ E_8$.~\qed\end{proof}
  \step{s3c5Amap}{An $\typ E_8$ automorphism carries this inclusion into $\typ E_7$.}
  \begin{proof}\pf\ Use the root assertion of Lemma~\ref{lem:sec3-choices}
   to send $\Zh$ to $\Zr_0$.~\qed\end{proof}
  \step{s3c5Aproper}{The inclusion has rank 7 on both sides and is proper.}
  \begin{proof}\pf\ Lemma~\ref{lem:sec3-spans} gives the ranks;
   Lemma~\ref{lem:3.4} gives $28<63$.~\qed\end{proof}
  \qedstep
  \begin{proof}\pf\ Combine \stepref{s3c5Amap} and \stepref{s3c5Aproper}.~\qed\end{proof}
 \end{proof}
 \step{s3c5AD}{\repair\ \assume{$\ell\geq4$ and an isometric inclusion $\typ A_\ell\subset \typ D_\ell$ exists.}
  \prove{FALSE.}}
 \begin{proof}
  \step{s3c5ADinduced}{The target friendship graph with $\ell-3$ triangles contains a coclique of size $\ell-2$.}
  \begin{proof}\pf\ Mark the source by its C4 star and apply
   Lemma~\ref{lem:sec3-marked-inclusion} and \stepref{s3c5uniform}.~\qed\end{proof}
  \step{s3c5friendbound}{A friendship graph with $t\geq1$ triangles has no coclique larger than $t$.}
  \begin{proof}
   \step{s3c5friendcentre}{A coclique containing the common vertex has size 1.}
   \begin{proof}\pf\ Every other vertex is adjacent to that vertex.~\qed\end{proof}
   \step{s3c5friendouter}{A coclique avoiding the common vertex has at most one vertex from each of the $t$ outer pairs.}
   \begin{proof}\pf\ Each outer pair is an edge.~\qed\end{proof}
   \qedstep
   \begin{proof}\pf\ These two cases give the bound $t$, since $t\geq1$.~\qed\end{proof}
  \end{proof}
  \qedstep
  \begin{proof}\pf\ Steps \stepref{s3c5ADinduced} and \stepref{s3c5friendbound}
   would give $\ell-2\leq\ell-3$.~\qed\end{proof}
 \end{proof}
 \step{s3c5Asix}{\assume{An isometric inclusion $\typ A_6\subset \typ E_6$ exists.}\prove{FALSE.}}
 \begin{proof}
  \step{s3c5gridfour}{The $3\times3$ lattice graph contains a 4-coclique.}
  \begin{proof}\pf\ C4 gives a 4-coclique for $\typ A_6$; use
   Lemma~\ref{lem:sec3-marked-inclusion} and \stepref{s3c5uniform}.~\qed\end{proof}
  \qedstep
  \begin{proof}\pf\ Among four grid positions two share one of the three rows,
   so those two are adjacent. This contradicts \stepref{s3c5gridfour}.~\qed\end{proof}
 \end{proof}
 \step{s3c5Dsix}{\assume{An isometric inclusion $\typ D_6\subset \typ E_6$ exists.}\prove{FALSE.}}
 \begin{proof}\pf\ The centre in the source marked graph has six neighbours by C4.
  Its image has at least six distinct neighbours by Lemma~\ref{lem:sec3-marked-inclusion};
  C4 gives target degree 4.~\qed\end{proof}
 \step{s3c5Dseven}{\assume{An isometric inclusion $\typ D_7\subset \typ E_7$ exists.}\prove{FALSE.}}
 \begin{proof}\pf\ The source centre has eight neighbours by C4; the induced
  injection of Lemma~\ref{lem:sec3-marked-inclusion} contradicts target degree 6.~\qed\end{proof}
 \qedstep
 \begin{proof}\pf\ The three positive inclusions and four obstructions give the assertion.~\qed\end{proof}
\end{proof}
\begingroup\raggedright
The coordinate inclusions already have Lean declarations
\leancite{LineSystem.An_Φ_subset_E₈viaA₈}{\texttt{LineSystem.\allowbreak An\_\allowbreak \ensuremath{\Phi}\_\allowbreak subset\_\allowbreak E\ensuremath{{}_8}viaA\ensuremath{{}_8}}} and \leancite{LineSystem.Dn_Φ_subset_E₈}{\texttt{LineSystem.\allowbreak Dn\_\allowbreak \ensuremath{\Phi}\_\allowbreak subset\_\allowbreak E\ensuremath{{}_8}}}.
C5 is proved in Lean, in \path{LineSystem/Inclusions.lean},
\path{Computations/C5Inclusions.lean}, \path{Computations/C5Dn.lean} and
\path{Computations/C5.lean}. An isometric inclusion is
\path{LineSystem.IsometricInclusion}; the lemma on an equal-rank proper inclusion is
\path{LineSystem.not_isMaximal_of_proper_isometric_inclusion}, and the marked-inclusion lemma
is \path{LineSystem.Hyp24.orthGraph_embedding_of_subset}. Lean proves
$\typ A_7\subset\typ E_7$ and $\typ A_8\subset\typ E_8$ by explicit integer coordinate rows,
\leancite{LineSystem.An.inclusionE₇}{\texttt{LineSystem.\allowbreak An.\allowbreak inclusionE\ensuremath{{}_7}}} and \leancite{LineSystem.An.inclusionE₈}{\texttt{LineSystem.\allowbreak An.\allowbreak inclusionE\ensuremath{{}_8}}}, not
through C3; $\typ D_8\subset\typ E_8$ is \leancite{LineSystem.Dn.inclusionE₈}{\texttt{LineSystem.\allowbreak Dn.\allowbreak inclusionE\ensuremath{{}_8}}}. Step
\texttt{s3c5Dtransitive} is \leancite{LineSystem.Dn.exists_isometry_star}{\texttt{LineSystem.\allowbreak Dn.\allowbreak exists\_\allowbreak isometry\_\allowbreak star}}, with
\leancite{LineSystem.Hyp24.nonempty_orthGraph_iso_of_S_eq_Dn}{\texttt{LineSystem.\allowbreak Hyp24.\allowbreak nonempty\_\allowbreak orthGraph\_\allowbreak iso\_\allowbreak of\_\allowbreak S\_\allowbreak eq\_\allowbreak Dn}}. The four obstructions are
\leancite{LineSystem.not_embeds_An_Dn}{\texttt{LineSystem.\allowbreak not\_\allowbreak embeds\_\allowbreak An\_\allowbreak Dn}} (for $\ell\geq4$), \leancite{LineSystem.not_embeds_An_E₆}{\texttt{LineSystem.\allowbreak not\_\allowbreak embeds\_\allowbreak An\_\allowbreak E\ensuremath{{}_6}}},
\leancite{LineSystem.not_embeds_Dn_E₆}{\texttt{LineSystem.\allowbreak not\_\allowbreak embeds\_\allowbreak Dn\_\allowbreak E\ensuremath{{}_6}}} and \leancite{LineSystem.not_embeds_Dn_E₇}{\texttt{LineSystem.\allowbreak not\_\allowbreak embeds\_\allowbreak Dn\_\allowbreak E\ensuremath{{}_7}}}, and
\leancite{Computations.C5.inclusions_and_obstructions}{\texttt{Computations.\allowbreak C5.\allowbreak inclusions\_\allowbreak and\_\allowbreak obstructions}} bundles the three proper
inclusions with them. There is no finite search: the grid bound of step \texttt{s3c5Asix} is the
pigeonhole argument above, and the friendship bound and the two degree contradictions are
proofs above as well.\par
\endgroup

\begin{remark}
 \repair\ The paper derives $\typ A_7\subset\typ E_7$ from the statement that a 5-coclique
 embeds in the triangular graph $T(6)$ (Error 5). With the orthogonality convention of C4,
 the graph of $\typ A_7$ is a 5-coclique and the graph of $\typ E_7$ is the disjointness
 graph $\comp{T(6)}$; the witness is a 5-coclique of $\comp{T(6)}$, equivalently a 5-clique
 of $T(6)$, and the five pairs $\{1,j\}$ with $2\leq j\leq6$ pairwise intersect. We record
 the correction only; C5 does not use this witness, since it proves the inclusion by an
 explicit coordinate inclusion.
\end{remark}

\begin{cgss}{Theorem}{3.6}\label{thm:3.6}
 \lean{LineSystem.classification_of_isMaximal}
 \uses{def:2.2,lem:2.3,thm:3.5,lem:3.4,comp:C4,comp:C5,lem:sec3-spans,lem:sec3-transport,lem:sec3-extension,lem:sec3-proper-inclusion}
 \assume{$S$ is a nonempty indecomposable line system in a finite-dimensional real
 inner product space, and $W=\operatorname{span}\Phi(S)$.}
 \prove{$S$ is maximal within $W$ if and only if it is isomorphic to one of
 \[
 \typ A_\ell\ (\ell\geq1,\ \ell\ne7,8),\qquad
 \typ D_\ell\ (\ell\geq4,\ \ell\ne8),\qquad \typ E_6,\ \typ E_7,\ \typ E_8.
 \]
 \repair\ Maximality is within the span (Implicit hypothesis 7), including for
 $\typ A_\ell\subset\re^{\ell+1}$.}
\end{cgss}
\begin{proof}
 \pfsketch\ A maximal system is star-closed. For the converse we will put any
 hypothetical proper extension inside an indecomposable star-closed system of the
 same rank, then exhaust its type using counts and C5.
 \step{s336forward}{\assume{$S$ is maximal within $W$.}
  \prove{$S$ has a type in the displayed list.}}
 \begin{proof}
  \step{s336closed}{$S$ is star-closed.}
  \begin{proof}\pf\ Apply Lemma~\ref{lem:2.3}.~\qed\end{proof}
  \step{s336classified}{$S$ has a type from Theorem 3.5.}
  \begin{proof}\pf\ Apply Theorem~\ref{thm:3.5} to \stepref{s336closed} and the standing hypotheses.~\qed\end{proof}
  \step{s336excluded}{Its type is not $\typ A_7$, $\typ A_8$, or $\typ D_8$.}
  \begin{proof}
   \step{s336exclAseven}{$\typ A_7$ is not maximal within its span.}
   \begin{proof}\pf\ Apply Lemma~\ref{lem:sec3-proper-inclusion} to the proper
    isometric inclusion $\typ A_7\subset\typ E_7$ of equal rank in C5.~\qed\end{proof}
   \step{s336exclAeight}{$\typ A_8$ is not maximal within its span.}
   \begin{proof}\pf\ Apply Lemma~\ref{lem:sec3-proper-inclusion} to the proper
    isometric inclusion $\typ A_8\subset\typ E_8$ of equal rank in C5.~\qed\end{proof}
   \step{s336exclDeight}{$\typ D_8$ is not maximal within its span.}
   \begin{proof}\pf\ Apply Lemma~\ref{lem:sec3-proper-inclusion} to the proper
    isometric inclusion $\typ D_8\subset\typ E_8$ of equal rank in C5.~\qed\end{proof}
   \qedstep
   \begin{proof}\pf\ A span isometry from $S$ to one of these models would carry
    maximality within $W$, from \asm{1}, to the model, by
    Lemma~\ref{lem:sec3-transport}; this contradicts \stepref{s336exclAseven},
    \stepref{s336exclAeight} or \stepref{s336exclDeight}.~\qed\end{proof}
  \end{proof}
  \qedstep
  \begin{proof}\pf\ Remove the three types of \stepref{s336excluded} from
   \stepref{s336classified}.~\qed\end{proof}
 \end{proof}
 \step{s336modelsmaximal}{Every retained model is maximal in its span.}
 \begin{proof}
  \step{s336contra}{\assume{$M$ is a retained model and is not maximal in its span.}
   \prove{FALSE.}}
  \begin{proof}
   \step{s336extension}{There is an indecomposable star-closed proper extension $T$ with the same rank as $M$.}
   \begin{proof}\pf\ C4 gives that $M$ is nonempty and indecomposable;
    apply Lemma~\ref{lem:sec3-extension} to $M$ in its span.~\qed\end{proof}
   \step{s336target}{There is a proper isometric inclusion of $M$ in a listed model
    of Theorem 3.5 of the same rank, with strictly more lines.}
   \begin{proof}
    \pflet{$W_M=\operatorname{span}\Phi(M)$.}
    \step{s336targethyp}{$T$ is a nonempty indecomposable star-closed line system in
     $W_M$, and $\operatorname{span}\Phi(T)=W_M$.}
    \begin{proof}\pf\ Lemma~\ref{lem:sec3-extension} gives these properties with the
     ambient space $W_M$, and $\Phi(M)\subseteq\Phi(T)$ makes $T$ nonempty.~\qed\end{proof}
    \step{s336targetiso}{There are a listed model $Y$ and a surjective linear isometry
     $G:W_M\to\operatorname{span}\Phi(Y)$ with $G(\Phi(T))=\Phi(Y)$.}
    \begin{proof}\pf\ Apply the forward direction of Theorem~\ref{thm:3.5} to $T$ in the
     finite-dimensional space $W_M$, using \stepref{s336targethyp}; its span is
     $W_M$.~\qed\end{proof}
    \step{s336targetmap}{$G$ is an isometric inclusion of $M$ in $Y$ with
     $G(\Phi(M))\subsetneq\Phi(Y)$.}
    \begin{proof}\pf\ $G$ is defined on $W_M=\operatorname{span}\Phi(M)$, it is injective,
     and $\Phi(M)\subsetneq\Phi(T)$ by \stepref{s336extension}; apply
     \stepref{s336targetiso}.~\qed\end{proof}
    \step{s336targetrank}{$Y$ has the rank of $M$.}
    \begin{proof}\pf\ $G$ is a linear bijection $W_M\to\operatorname{span}\Phi(Y)$ by
     \stepref{s336targetiso}.~\qed\end{proof}
    \qedstep
    \begin{proof}\pf\ A proper inclusion of finite antipodal sets strictly increases
     the number of pairs; use \stepref{s336targetmap} and \stepref{s336targetrank}.~\qed\end{proof}
   \end{proof}
   \step{s336E}{\case{$M=\typ E_r$ for $r\in\{6,7,8\}$.}}
   \begin{proof}\pf\ At those ranks the $(\typ A,\typ D,\typ E)$ line counts in
    Lemma~\ref{lem:3.4} are $(21,30,36)$, $(28,42,63)$, and $(36,56,120)$.
    The $\typ E$ entry is largest in each row, contrary to \stepref{s336target}.~\qed\end{proof}
   \step{s336A}{\case{$M=\typ A_\ell$, $\ell\ne7,8$.}}
   \begin{proof}
    \step{s336Asmall}{\case{$\ell\leq3$.}}
    \begin{proof}\pf\ Theorem~\ref{thm:3.5} has only $\typ A_\ell$ at this rank,
     so the strict count increase in \stepref{s336target} is impossible.~\qed\end{proof}
    \step{s336Alarge}{\case{$\ell\geq4$.}}
    \begin{proof}
     \step{s336AD}{The target cannot be $\typ D_\ell$.}
     \begin{proof}\pf\ C5 excludes $\typ A_\ell\subset \typ D_\ell$ for every $\ell\geq4$.~\qed\end{proof}
     \step{s336AE}{The target cannot be exceptional.}
     \begin{proof}\pf\ Rank equality and $\ell\ne7,8$ leave only $\typ E_6$;
      C5 excludes $\typ A_6\subset \typ E_6$.~\qed\end{proof}
     \qedstep
     \begin{proof}\pf\ The only remaining target is $\typ A_\ell$, with equal count,
      contrary to \stepref{s336target}.~\qed\end{proof}
    \end{proof}
    \qedstep
    \begin{proof}\pf\ Steps \stepref{s336Asmall} and \stepref{s336Alarge} exhaust positive ranks.~\qed\end{proof}
   \end{proof}
   \step{s336D}{\case{$M=\typ D_\ell$, $\ell\geq4$, $\ell\ne8$.}}
   \begin{proof}
    \step{s336DA}{The target cannot be $\typ A_\ell$ or $\typ D_\ell$.}
    \begin{proof}\pf\ Their counts are at most that of $\typ D_\ell$, since
     $\ell(\ell-1)-\ell(\ell+1)/2=\ell(\ell-3)/2>0$;
     this contradicts \stepref{s336target}.~\qed\end{proof}
    \step{s336DE}{The target cannot be exceptional.}
    \begin{proof}\pf\ With equal ranks and $\ell\ne8$, only $\typ D_6\subset \typ E_6$
     and $\typ D_7\subset \typ E_7$ remain; C5 excludes both.~\qed\end{proof}
    \qedstep
    \begin{proof}\pf\ Theorem~\ref{thm:3.5} has no further target types.~\qed\end{proof}
   \end{proof}
   \qedstep
   \begin{proof}\pf\ The retained model list is exhausted by
    \stepref{s336E}, \stepref{s336A}, and \stepref{s336D}.~\qed\end{proof}
  \end{proof}
  \qedstep
  \begin{proof}\pf\ Step \stepref{s336contra} rules out nonmaximality of any retained model.~\qed\end{proof}
 \end{proof}
 \step{s336reverse}{\assume{$S$ is isomorphic to a retained model.}
  \prove{$S$ is maximal within $W$.}}
 \begin{proof}\pf\ Transport \stepref{s336modelsmaximal} by
  Lemma~\ref{lem:sec3-transport}.~\qed\end{proof}
 \qedstep
 \begin{proof}\pf\ Steps \stepref{s336forward} and \stepref{s336reverse} give the equivalence.~\qed\end{proof}
\end{proof}

\begin{namedresult}{Definition}{(signings and the matrix convention)}\label{def:sec3-signing}
 \lean{LineSystem.Signing}\lean{LineSystem.signedAdj}\lean{LineSystem.nonempty_signing}
 \uses{def:line-system,facts:inner}
 \assume{$S$ is a line system, $W=\operatorname{span}\Phi(S)$, and $I$ is a finite index set.}
 \new\ We make ``one vector per line'' an explicit bijection.
 \pflet{A signing is a family $\sigma:I\to\Phi(S)$ for which
 $(i,\epsilon)\mapsto\epsilon\sigma(i)$ is a bijection
 $I\times\{1,-1\}\to\Phi(S)$. Put $N=|I|$,
 $M(i,j)=(\sigma(i),\sigma(j))$, and $C=M-2I_N$.}
 Thus $N$ is the number of lines. The diagonal of $C$ is zero, and its off-diagonal
 entries lie in $\{0,1,-1\}$ by Section~\ref{facts:inner}. A change of signs conjugates
 $C$ by a diagonal matrix of signs, and a reordering conjugates by a permutation matrix.
 We use $I_N$ for the identity matrix, to distinguish it from the index set $I$.
 Every line system has a signing indexed by $\{1,\dots,N\}$, where $N$ is its number of lines:
 choose one of the two vectors on each line (\path{LineSystem.nonempty_signing}). Thus the
 equations of Theorem 3.7, which hold for every signing, are not vacuous.
\end{namedresult}

\begin{namedresult}{Computation}{C6}\label{comp:C6}
 \lean{Computations.C6.coordinate_frame_matrices}
 \uses{def:3.1,def:3.2,def:3.3,lem:sec3-integer,comp:C2}
 \assume{$R$ is the complete, duplicate-free list of raw signed roots from
 Lemma~\ref{lem:sec3-integer}, with the coordinate cuts for $\typ E_7,\typ E_6$ from
 Definition~\ref{def:3.3}. Write
 $K=\sum_{v\in R}vv^T$ and let $J_m$ be the $m\times m$ all-one matrix.}
 \prove{The following are exact integer matrix identities:
 \[
 \begin{array}{c|c|c}
 \text{model}&q&K\\\hline
 \typ A_\ell&1&2((\ell+1)I_{\ell+1}-J_{\ell+1})\\
 \typ D_\ell&1&4(\ell-1)I_\ell\\
 \typ E_8\text{ (i)}&9&540I_9-60J_9\\
 \typ E_8\text{ (ii)}&4&240I_8\\
 \typ E_8\text{ (iii)}&2&120I_8\\
 \typ E_7&4&\operatorname{diag}(144I_6,72J_2)\\
 \typ E_6&4&\operatorname{diag}(96I_5,32J_3).
 \end{array}
 \]}
\end{namedresult}
\begin{proof}
 \pfsketch\ \new\ We replace spectral calculations by small coordinate frame
 matrices. The classical rows are parametric; the exceptional rows use the
 existing Gram checker on the columns of the raw root table.
 \step{s3c6A}{The $\typ A_\ell$ identity holds for every $\ell\geq1$.}
 \begin{proof}
  \step{s3c6Adiag}{$K(i,i)=2\ell$.}
  \begin{proof}\pf\ There are $\ell$ difference roots with entry 1 at $i$ and
   $\ell$ with entry $-1$ there.~\qed\end{proof}
  \step{s3c6Aoff}{For $i\ne j$, $K(i,j)=-2$.}
  \begin{proof}\pf\ Only $\Ze_i-\Ze_j$ and its negative contribute; both products are $-1$.~\qed\end{proof}
  \qedstep
  \begin{proof}\pf\ These entries are those of $2((\ell+1)I-J)$.~\qed\end{proof}
 \end{proof}
 \step{s3c6D}{The $\typ D_\ell$ identity holds for every $\ell\geq4$.}
 \begin{proof}
  \step{s3c6Ddiag}{$K(i,i)=4(\ell-1)$.}
  \begin{proof}\pf\ Each support $\{i,j\}$ with $j\ne i$ has four sign choices, each contributing 1.~\qed\end{proof}
  \step{s3c6Doff}{For $i\ne j$, $K(i,j)=0$.}
  \begin{proof}\pf\ The four signs on the only contributing support $\{i,j\}$
   give products $1,-1,-1,1$.~\qed\end{proof}
  \qedstep
  \begin{proof}\pf\ The entries in the two steps give $4(\ell-1)I$.~\qed\end{proof}
 \end{proof}
 \step{s3c6exceptional}{The five exceptional raw identities hold.}
 \begin{proof}
  \pflet{$P_i=(v_i)_{v\in R}$ be the $i$th column of the raw root table, in its
   list order. The list $P$ has length equal to the ambient dimension.}
  \step{s3c6one}{For model (i), \texttt{checkGram} accepts $P$ and $540I_9-60J_9$.}
  \begin{proof}\pf\ C6 checks \path{Computations.checkGram} with vector
   dimension 240 and nine columns, from the pair and signed-triple list of C2.~\qed\end{proof}
  \step{s3c6two}{For model (ii), \texttt{checkGram} accepts $P$ and $240I_8$.}
  \begin{proof}\pf\ C6 applies that checker with dimension 240 and eight columns
   of the doubled pair and even-sign list.~\qed\end{proof}
  \step{s3c6three}{For model (iii), \texttt{checkGram} accepts $P$ and $120I_8$.}
  \begin{proof}\pf\ C6 applies that checker with dimension 240 and eight columns
   of the axis and block-sign list.~\qed\end{proof}
  \step{s3c6seven}{For $\typ E_7$, \texttt{checkGram} accepts $P$ and $\operatorname{diag}(144I_6,72J_2)$.}
  \begin{proof}\pf\ C6 applies that checker with dimension 126 and eight columns
   of the doubled list filtered by $v_7=v_8$.~\qed\end{proof}
  \step{s3c6six}{For $\typ E_6$, \texttt{checkGram} accepts $P$ and $\operatorname{diag}(96I_5,32J_3)$.}
  \begin{proof}\pf\ C6 applies that checker with dimension 72 and eight columns
   of the doubled list filtered by $v_6=v_7=v_8$.~\qed\end{proof}
  \qedstep
  \begin{proof}\pf\ For each accepted table,
   \path{Computations.checkGram_sound} gives
   $K(i,j)=P_i\cdot P_j=\sum_{v\in R}v_iv_j$, including diagonal entries.~\qed\end{proof}
 \end{proof}
 \qedstep
 \begin{proof}\pf\ The two parametric cases and five column-Gram checks cover
  the seven rows.~\qed\end{proof}
\end{proof}

\begingroup\raggedright
C6 is proved in Lean, in \path{LineSystem/Frame.lean} and \path{Computations/C6Classical.lean},
\path{Computations/C6E8.lean}, \path{Computations/C6E7E6.lean} and \path{Computations/C6.lean}.
Lean states C6 as the frame property on the span, \path{LineSystem.IsFrame}, equivalent to the
column-Gram form by \path{LineSystem.isFrame_iff_gram} and transported by
\path{LineSystem.Iso.isFrame}. The two parametric rows are \path{Computations.C6.an_frame_matrix}
and \path{Computations.C6.dn_frame_matrix}, giving \path{LineSystem.isFrame_An} and
\path{LineSystem.isFrame_Dn}; the exceptional matrices are checked by \texttt{decide +kernel} over
structured indices, not by \path{Computations.checkGram}, and give
\leancite{LineSystem.isFrame_E₈}{\texttt{LineSystem.\allowbreak isFrame\_\allowbreak E\ensuremath{{}_8}}}, \leancite{LineSystem.isFrame_E₇}{\texttt{LineSystem.\allowbreak isFrame\_\allowbreak E\ensuremath{{}_7}}} and \leancite{LineSystem.isFrame_E₆}{\texttt{LineSystem.\allowbreak isFrame\_\allowbreak E\ensuremath{{}_6}}}.
The bundle is \path{Computations.C6.coordinate_frame_matrices}. The checker
\path{Computations.checkGram} is defined in \path{Computations/ListData.lean}, and its
soundness theorem \path{Computations.checkGram_sound} is in \path{Computations/Checkers.lean};
C6 does not use them. The frame lemma and Theorem~3.7 below are proved in Lean as well, the
frame lemma with the rank of the model (\path{LineSystem.frame_identity}) and with the rank of
the span (\path{LineSystem.frame_identity_span}). The three span lemmas are in Lean:
\texttt{LineSystem.\allowbreak Iso.\allowbreak spanIsometry} with the transport of
the predicates and of the marked sets, \texttt{LineSystem.\allowbreak restrictSpan} (with
\texttt{LineSystem.\allowbreak Iso.\allowbreak restrictSpan} and \texttt{LineSystem.\allowbreak Iso.\allowbreak ofSpanIsometry},
the interface of the final step of Theorem~3.5) in \path{LineSystem/SpanTransport.lean}, and
\texttt{LineSystem.\allowbreak exists\_\allowbreak maximal\_\allowbreak extension\_\allowbreak span} in
\path{LineSystem/SpanExtension.lean}; \path{LineSystem/SpanAudit.lean} prints their axioms. The Lean
proofs of C5 and C6 transport along \texttt{LineSystem.Iso} and use the model's rank.\par
Theorems~3.5 and~3.6 are in Lean. Theorem~3.5 is
\texttt{LineSystem.\allowbreak classification} in \path{LineSystem/Classification.lean}, with its
forward direction \texttt{LineSystem.\allowbreak exists\_\allowbreak iso\_\allowbreak rootType};
Theorem~3.6 is \texttt{LineSystem.\allowbreak classification\_\allowbreak of\_\allowbreak isMaximal}
in \path{LineSystem/ClassificationMaximal.lean}, with the form that Section~4 uses,
\texttt{LineSystem.\allowbreak IsMaximal.\allowbreak exists\_\allowbreak iso\_\allowbreak rootType}.
Both rest on the unconditional Theorem~2.13,
\texttt{LineSystem.\allowbreak five\_\allowbreak possibilities'}, whose inputs from Section~3 are
the packaged facts \texttt{LineSystem.\allowbreak modelInputs}, proved from C2 and C4
without either classification, in \path{LineSystem/ClassificationModels.lean}.
Lemma (a connected spanning set) is
\texttt{LineSystem.\allowbreak isIndecomposable\_\allowbreak of\_\allowbreak connected\_\allowbreak spanning}
in \path{LineSystem/ClassificationConnected.lean}, and the perpendicular choices of
Lemma (independence of the perpendicular choices) in the default model are
\leancite{LineSystem.E₈.nonempty_iso_E₇_orthogonalTo}{\texttt{LineSystem.\allowbreak E\ensuremath{{}_8}.\allowbreak nonempty\_\allowbreak iso\_\allowbreak E\ensuremath{{}_7}\_\allowbreak orthogonalTo}},
\leancite{LineSystem.E₈.nonempty_iso_E₆_orthogonalTo}{\texttt{LineSystem.\allowbreak E\ensuremath{{}_8}.\allowbreak nonempty\_\allowbreak iso\_\allowbreak E\ensuremath{{}_6}\_\allowbreak orthogonalTo}}
and
\leancite{LineSystem.E₈.exists_isometry_of_ordered_star}{\texttt{LineSystem.\allowbreak E\ensuremath{{}_8}.\allowbreak exists\_\allowbreak isometry\_\allowbreak of\_\allowbreak ordered\_\allowbreak star}}
in \path{LineSystem/ClassificationChoices.lean}. The Lean departs from the proofs above in three
places. The ambient space of Theorems~3.5 and~3.6 need not be finite-dimensional, since the
span of the finite root set is, and Theorem~2.13 is applied there. The forward direction of
Theorem~3.5 is shorter: the isomorphisms of Theorem~2.13 go directly to the named models and are
composed with \texttt{LineSystem.\allowbreak Iso.\allowbreak restrictSpan}, so steps
\texttt{s335coords} and \texttt{s335orthogonal} of the proof below are not needed. The retained list of
Theorem~3.6 is stated as an explicit disjunction of the five families.\par
The two remaining pieces are in Lean as well, and their modules are imported by
\path{LineGraphs.lean}. The bundles
\texttt{Computations.\allowbreak C2.\allowbreak model\_\allowbreak presentations} and
\texttt{Computations.\allowbreak C2.\allowbreak line\_\allowbreak counts} of
Lemma (integer realizations) and C2 are in \path{Computations/C2.lean}, with the cuts
\leancite{LineSystem.T₇}{\texttt{LineSystem.\allowbreak T\ensuremath{{}_7}}} and
\leancite{LineSystem.T₆}{\texttt{LineSystem.\allowbreak T\ensuremath{{}_6}}} of
Definition~3.3; Lean checks their 126 and 72 roots in the kernel, as for $\typ E_7$
and $\typ E_6$. The assertions of Lemma (independence of the perpendicular choices) for the
other two $\typ E_8$ models, with $T_7\cong\typ E_7$ and $T_6\cong\typ E_6$, are in
\path{LineSystem/ClassificationChoicesModels.lean}. In the final step the Lean departs from the
proof above: a C3 isomorphism is a bijection of roots, not a map of the ambient space, so Lean
carries the automorphisms of the default model across it through the span isometry of
Lemma (transport on spans), extended to the whole space (for model (i), from the zero-sum
hyperplane to $\re^9$). The file \path{Computations/C2ChoicesAudit.lean} prints the axioms of
all these declarations.\par
\endgroup

\begin{namedresult}{Lemma}{(the frame identity within the span)}\label{lem:sec3-frame}
 \lean{LineSystem.frame_identity}
 \uses{comp:C6,comp:C2,lem:sec3-spans,def:sec3-signing,lem:sec3-transport,facts:inner}
 \assume{$S$ is one of the listed models, or an isometric copy of it;
 $W=\operatorname{span}\Phi(S)$ has rank $\ell>0$ and $N$ lines.
 Let $\sigma:I\to\Phi(S)$ be any signing.}
 \prove{For every $\Zx\in W$,
 \[
 \sum_{\Zr\in\Phi(S)}(\Zr,\Zx)\Zr=\frac{4N}{\ell}\Zx,
 \qquad
 \sum_{i\in I}(\sigma(i),\Zx)\sigma(i)=\frac{2N}{\ell}\Zx.
 \]
 \repair\ These formulas are asserted on $W$, not on a larger coordinate space
 (Implicit hypothesis 7 and the model-span correction).}
\end{namedresult}
\begin{proof}
 \pfsketch\ We will divide C6 by the raw scale $q$, restrict to the known span,
 and halve the sum using antipodal pairs.
 \step{s3frameoperator}{For a coordinate model, the full signed-root operator is $K/q$.}
 \begin{proof}\pf\ Substitute $\Zr=c\Zv$ in the sum, using $c^2=1/q$.
  The exact no-multiplicity transfer is \leancite{LineSystem.IntegerModel.sum_Φ}{\texttt{LineSystem.\allowbreak IntegerModel.\allowbreak sum\_\allowbreak \ensuremath{\Phi}}}.~\qed\end{proof}
 \step{s3framerestrict}{The operator is scalar on $W$, with scalar $4N/\ell$.}
 \begin{proof}
  \step{s3frameA}{On the $\typ A_\ell$ span, $K/q=2(\ell+1)I$.}
  \begin{proof}\pf\ $J\Zx=0$ on the zero-sum span in Lemma~\ref{lem:sec3-spans};
   use the first row of C6.~\qed\end{proof}
  \step{s3frameD}{On the $\typ D_\ell$ span, $K/q=4(\ell-1)I$.}
  \begin{proof}\pf\ Its C6 matrix has that value on the whole coordinate space.~\qed\end{proof}
  \step{s3frameEi}{On model (i)'s span, $K/q=60I$.}
  \begin{proof}\pf\ Its span has sum zero, so $(540I-60J)\Zx/9=60\Zx$.~\qed\end{proof}
  \step{s3frameEii}{On model (ii)'s span, $K/q=60I$.}
  \begin{proof}\pf\ Divide $240I_8$ by $q=4$.~\qed\end{proof}
  \step{s3frameEiii}{On model (iii)'s span, $K/q=60I$.}
  \begin{proof}\pf\ Divide $120I_8$ by $q=2$.~\qed\end{proof}
  \step{s3frameSeven}{On the $\typ E_7$ span, $K/q=36I$.}
  \begin{proof}\pf\ On a last-coordinate pair $(t,t)$, $72J_2(t,t)=144(t,t)$;
   the first block is $144I_6$. Divide by 4.~\qed\end{proof}
  \step{s3frameSix}{On the $\typ E_6$ span, $K/q=24I$.}
  \begin{proof}\pf\ On $(t,t,t)$, $32J_3(t,t,t)=96(t,t,t)$;
   the first block is $96I_5$. Divide by 4.~\qed\end{proof}
  \qedstep
  \begin{proof}\pf\ Substitution of the C2 line counts gives $4N/\ell$ in each row.~\qed\end{proof}
 \end{proof}
 \step{s3framecopies}{The full-root formula holds for isometric copies in their spans.}
 \begin{proof}\pf\ Conjugate the full-root operator through the span isometry
  of Lemma~\ref{lem:sec3-transport}; reindex by its root bijection.~\qed\end{proof}
 \step{s3framehalf}{The full-root sum equals twice the signing sum.}
 \begin{proof}
  \step{s3framepair}{The roots $\Zr,-\Zr$ contribute the same summand.}
  \begin{proof}\pf\ $(-\Zr,\Zx)(-\Zr)=(\Zr,\Zx)\Zr$ by bilinearity.~\qed\end{proof}
  \qedstep
  \begin{proof}\pf\ Reindex the full-root sum by the signing bijection in
   Definition~\ref{def:sec3-signing} and use \stepref{s3framepair}.~\qed\end{proof}
 \end{proof}
 \qedstep
 \begin{proof}\pf\ The full formula follows from \stepref{s3frameoperator}--\stepref{s3framecopies};
  divide the equality in \stepref{s3framehalf} by 2 for the signing formula.~\qed\end{proof}
\end{proof}

\begin{namedresult}{Lemma}{(the Gram polynomial from a frame)}\label{lem:sec3-gram-polynomial}
 \lean{LineSystem.signedAdj_quadratic_of_frame}
 \uses{def:sec3-signing,facts:inner}
 \assume{\begin{pfenum}
  \item $S$ is a line system and $W=\operatorname{span}\Phi(S)$.
  \item $\sigma:I\to\Phi(S)$ is a signing.
  \item $t\in\re$ satisfies $\sum_i(\sigma(i),\Zx)\sigma(i)=t\Zx$ for every $\Zx\in W$.
 \end{pfenum}}
 \prove{With $M=2I_N+C$ as in Definition~\ref{def:sec3-signing},
 $M^2=tM$ and $(C+2I_N)(C-(t-2)I_N)=0$.}
\end{namedresult}
\begin{proof}
 \pfsketch\ \new\ We isolate the algebraic passage from a frame to the signed
 Gram polynomial. We will apply the frame equation to each root and take its product
 with each other root. This avoids choosing an ambient orthonormal basis.
 \step{s3gramentry}{For all $i,j\in I$, $(M^2)(i,j)=tM(i,j)$.}
 \begin{proof}
  \step{s3gramsum}{$(M^2)(i,j)=(\sigma(i),\sum_k(\sigma(k),\sigma(j))\sigma(k))$.}
  \begin{proof}\pf\ Expand matrix multiplication and then bilinearity from Section~\ref{facts:inner}.~\qed\end{proof}
  \step{s3gramapply}{The vector sum in \stepref{s3gramsum} is $t\sigma(j)$.}
  \begin{proof}\pf\ Apply the assumed frame identity to $\sigma(j)\in\Phi(S)\subseteq W$.~\qed\end{proof}
  \qedstep
  \begin{proof}\pf\ Substitute \stepref{s3gramapply} into \stepref{s3gramsum} and use scalar linearity.~\qed\end{proof}
 \end{proof}
 \step{s3grammatrix}{$M(M-tI_N)=0$.}
 \begin{proof}\pf\ Step \stepref{s3gramentry} gives $M^2=tM$ entry by entry;
  distribute multiplication by $M-tI_N$.~\qed\end{proof}
 \qedstep
 \begin{proof}\pf\ Substitute $M=C+2I_N$ in \stepref{s3grammatrix};
  then $M-tI_N=C-(t-2)I_N$.~\qed\end{proof}
\end{proof}

\begin{cgss}{Theorem}{3.7}\label{thm:3.7}
 \lean{LineSystem.signedAdj_quadratic_An}\lean{LineSystem.signedAdj_quadratic_Dn}
 \lean{LineSystem.signedAdj_quadratic_E₈}\lean{LineSystem.signedAdj_quadratic_E₇}
 \lean{LineSystem.signedAdj_quadratic_E₆}
 \uses{def:sec3-signing,lem:3.4,lem:sec3-spans,comp:C6,lem:sec3-frame,lem:sec3-gram-polynomial}
 \assume{$S$ is $\typ A_\ell$ with $\ell\geq1$, $\typ D_\ell$ with $\ell\geq4$, or one
 of $\typ E_8,\typ E_7,\typ E_6$, in any isometric realization. Choose one vector of squared
 norm 2 on each line, and let $2I_N+C$ be their Gram matrix.}
 \prove{For every choice and ordering of those vectors,
 \[
 \begin{array}{rcl}
 \typ A_\ell:&(C+2I_N)(C-(\ell-1)I_N)&=0,\\
 \typ D_\ell:&(C+2I_N)(C-2(\ell-2)I_N)&=0,\\
 \typ E_8:&(C+2I_N)(C-28I_N)&=0,\\
 \typ E_7:&(C+2I_N)(C-16I_N)&=0,\\
 \typ E_6:&(C+2I_N)(C-10I_N)&=0.
 \end{array}
 \]
 \repair\ The $\typ E_8$ constant is 28, correcting the printed 18 (Error 1).}
\end{cgss}
\begin{proof}
 \pfsketch\ We will use C6's frame identity within the span. The second factor
 has constant $\theta=2N/\ell-2$, with $\ell$ the rank, not the number of ambient coordinates.
 \step{s337general}{$(C+2I_N)(C-(2N/\ell-2)I_N)=0$.}
 \begin{proof}\pf\ Apply Lemma~\ref{lem:sec3-gram-polynomial} to the signing
  formula of Lemma~\ref{lem:sec3-frame}.~\qed\end{proof}
 \step{s337A}{For $\typ A_\ell$, $2N/\ell-2=\ell-1$.}
 \begin{proof}\pf\ Lemma~\ref{lem:3.4} gives $N=\ell(\ell+1)/2$;
  cancel the positive $\ell$.~\qed\end{proof}
 \step{s337D}{For $\typ D_\ell$, $2N/\ell-2=2(\ell-2)$.}
 \begin{proof}\pf\ Substitute $N=\ell(\ell-1)$ and cancel $\ell>0$.~\qed\end{proof}
 \step{s337Eeight}{For $\typ E_8$, $2N/\ell-2=28$.}
 \begin{proof}\pf\ The rank and count are 8 and 120 by
  Lemmas~\ref{lem:sec3-spans} and~\ref{lem:3.4}; $240/8-2=28$.~\qed\end{proof}
 \step{s337Eseven}{For $\typ E_7$, $2N/\ell-2=16$.}
 \begin{proof}\pf\ Its rank and count are 7 and 63, so $126/7-2=16$.~\qed\end{proof}
 \step{s337Esix}{For $\typ E_6$, $2N/\ell-2=10$.}
 \begin{proof}\pf\ Its rank and count are 6 and 36, so $72/6-2=10$.~\qed\end{proof}
 \qedstep
 \begin{proof}\pf\ Substitute the five constants into \stepref{s337general}.
  The frame lemma applies to every signing and every isometric copy, so no
  additional switching or eigenvalue calculation is needed.~\qed\end{proof}
\end{proof}
\endgroup
